% concatenated from part2.tex
\documentclass{article}
\usepackage{amsmath, amsthm, amssymb}
\usepackage{hyperref}
\hypersetup{hidelinks}
\usepackage{geometry}
\geometry{margin=1in}
\usepackage[most]{tcolorbox}
\tcbset{
  keybox/.style={
    enhanced,
    colback=gray!8,
    colframe=gray!50,
    boxrule=0.5pt,
    arc=2pt,
    left=6pt, right=6pt, top=6pt, bottom=6pt,
    fontupper=\small
  }
}

\newtheorem{theorem}{Theorem}[section]
\newtheorem{proposition}[theorem]{Proposition}
\newtheorem{lemma}[theorem]{Lemma}
\newtheorem{corollary}[theorem]{Corollary}
\newtheorem{definition}[theorem]{Definition}
\newtheorem{remark}[theorem]{Remark}

\title{Universal Shuffle Asymptotics, Part II:\\
Non-Gaussian Limits for Shuffle Privacy---Poisson, Skellam,\\
and Compound-Poisson Regimes}
\author{Alex Shvets}
\date{March 2026}

\begin{document}

\maketitle

\begin{abstract}
Part I of this series~\cite{Shv26} develops a sharp Gaussian (LAN/GDP) limit theory for neighboring
shuffle experiments when the local randomizer is fixed and has full support bounded away from
zero. The present paper characterizes the first universality-breaking frontier: critical sequences
of increasingly concentrated local randomizers for which classical Lindeberg conditions fail and
the shuffle score exhibits rare macroscopic jumps.

For shuffled binary randomized response with local privacy $\varepsilon_0 = \varepsilon_0(n)$, we prove experiment-level
convergence (in Le Cam distance) to explicit shift limit experiments: a Poisson-shift limit for
the canonical neighboring pair when $\exp(\varepsilon_0(n))/n \to c^2$, and a Skellam-shift limit for proportional
compositions $k/n \to \pi \in (0,1)$ in the same scaling, including an explicit disappearance of the
two-sided $\delta$-floor away from boundary compositions.

For general finite alphabets, we introduce a sparse-error critical regime and prove a multivariate
compound-Poisson / independent Poisson vector limit for the centered released histogram, yielding
a multivariate Poisson-shift experiment and an explicit limiting $(\varepsilon, \delta)$ curve as a multivariate
Poisson series. Together with Part~I, these results yield a three-regime picture
(Gaussian/GDP, critical Poisson/Skellam/compound-Poisson, and super-critical no privacy)
under convergent macroscopic scalings.
\end{abstract}

\noindent\textbf{MSC 2020:} 62B15 (statistical experiments and information); 68P27 (privacy); 60F05 (central limit and other weak theorems); 60E07 (infinitely divisible distributions).\\[2pt]
\noindent\textbf{Keywords:} shuffle model, differential privacy, Le Cam distance, Poisson approximation, Skellam distribution, compound Poisson, privacy amplification.

\tableofcontents

\section{Introduction}

The shuffle model can yield strong amplification of local differential privacy (LDP) by anonymizing
the multiset of local messages~\cite{EFM19,CSU19,BBG19}. Part I~\cite{Shv26} gives a sharp Gaussian characterization---including LAN/GDP equivalence in the sense of Gaussian differential privacy~\cite{DRS22}
and explicit privacy curves---when the local randomizer is fixed and has full support bounded away from 0.
Independently, Takagi and Liew~\cite{TL26} develop an asymptotic blanket-divergence framework for shuffle
privacy beyond pure LDP; the present paper is complementary in focusing on critical non-Gaussian limit
experiments.\footnote{Part~I of this series~\cite{Shv26} was submitted to arXiv on 17~January~2026
(v1, as recorded in the arXiv submission history). The identifier \texttt{2602.09029} reflects the
processing date rather than the submission date. Takagi--Liew~\cite{TL26} appeared on 27~January~2026.}

In practice, however, one often lets the local privacy level $\varepsilon_0 = \varepsilon_0(n)$ grow with the population
size $n$ (e.g.\ to reduce estimator variance). In such settings the shuffle score may fail to be a sum of
``small'' increments: extremely unlikely local outcomes can produce $\Theta(1)$ jumps in the log-likelihood
ratio at the population scale. The correct limit ceases to be Gaussian and becomes Poisson or
compound-Poisson.

This paper develops a sharp non-Gaussian limit theory at the critical threshold, at the level of binary
experiments and Le Cam convergence, with explicit total-variation rates and explicit limiting privacy curves.
Together with Part~I, it delineates how the critical and super-critical fronts sit around the Gaussian regime
whenever the relevant macroscopic composition and scaling parameters converge.

The main results are Theorems~3.1, 4.1, and~5.8 (Poisson-shift, Skellam-shift, and multivariate compound-Poisson limits, each with explicit $O(n^{-1})$ Le Cam rates and quantitative privacy curve convergence) together with the phase-diagram synthesis of Section~6 under convergent macroscopic scalings. Proposition~5.4 in Section~5.2 additionally develops a hybrid Gaussian/compound-Poisson weak limit for the two-dominant regime, characterizing the first L\'evy--Khintchine layer; Appendix~\ref{app:hybrid-privacy} closes the resulting privacy-curve convergence for interior compositions $\pi\in(0,1)$, while boundary compositions reduce to the Poisson-shift case of Section~3.

\section{Model and preliminaries}

\subsection{Shuffle mechanism and transcript laws}

Fix a population size $n \geq 1$. Each user $i \in \{1, \ldots, n\}$ holds a private datum $x_i \in \{0, 1\}$. A (possibly
$n$-dependent) local randomizer is a Markov kernel
\[
W^{(n)}: \{0, 1\} \to \Delta(\mathcal{Y}), \quad y \mapsto W^{(n)}_x(y),
\]
where $\mathcal{Y}$ is a finite output alphabet and $\Delta(\mathcal{Y})$ denotes the simplex of probability measures on $\mathcal{Y}$.
Given an input dataset $x^n = (x_1, \ldots, x_n)$, users apply $W^{(n)}$ independently to produce messages
$Y_i \sim W^{(n)}_{x_i}$.

The shuffle mechanism outputs the multiset of messages, equivalently the histogram
\[
N(y) := \sum_{i=1}^n \mathbf{1}\{Y_i = y\} \in \mathbb{Z}^{\mathcal{Y}}, \qquad \sum_{y \in \mathcal{Y}} N(y) = n.
\]

For $k \in \{0, \ldots, n\}$, let $T_{n,k}$ denote the shuffle transcript law when the dataset has exactly $k$ ones and
$n - k$ zeros. Since the shuffle output is permutation-invariant, $T_{n,k}$ depends on $x^n$ only through $k$.

\subsection{Privacy loss, privacy curves, and trade-off}

Consider two neighboring datasets differing in one entry, i.e.\ $k$ versus $k+1$ ones. The corresponding
neighboring shuffle experiment is the binary experiment
\[
\mathcal{E}_{n,k} := (P_{n,k}, Q_{n,k}), \quad P_{n,k} := T_{n,k}, \quad Q_{n,k} := T_{n,k+1}.
\]

Whenever $Q_{n,k} \ll P_{n,k}$, define the likelihood ratio and privacy loss random variables
\[
L_{n,k} := \frac{dQ_{n,k}}{dP_{n,k}}, \qquad \Lambda_{n,k} := \log L_{n,k}.
\]

\textbf{One-sided privacy curve.} For $\varepsilon \geq 0$, the tight one-sided privacy curve (privacy profile) is
\begin{equation}
\delta_{Q\|P}(\varepsilon) := \sup_A \{Q(A) - e^\varepsilon P(A)\}, \label{eq:delta}
\end{equation}
where the supremum is over measurable sets $A$.

\begin{lemma}[Neyman--Pearson identities for the privacy curve]
Assume $Q \ll P$ and let $L = dQ/dP$. Then for every $\varepsilon \geq 0$,
\begin{equation}
\delta_{Q\|P}(\varepsilon) = E_P[(L - e^\varepsilon)_+], \quad (u)_+ := \max\{u, 0\}. \label{eq:NP1}
\end{equation}
If the sample space is countable with pmfs $p$ and $q$, then for every $\varepsilon \geq 0$,
\begin{equation}
\delta_{Q\|P}(\varepsilon) = \sum_x \bigl(q(x) - e^\varepsilon p(x)\bigr)_+. \label{eq:NP2}
\end{equation}
The first identity requires $Q \ll P$; the second does not.
\end{lemma}

\begin{proof}
Under $Q \ll P$, for any measurable $A$ we may write
\[
Q(A) - e^\varepsilon P(A) = \int_A (L - e^\varepsilon)\, dP.
\]
The integrand is pointwise maximized by taking $A^* = \{L > e^\varepsilon\}$, hence
\[
\sup_A \int_A (L - e^\varepsilon)\, dP = \int (L - e^\varepsilon)_+\, dP,
\]
which is \eqref{eq:NP1}. On a countable space, for any $A \subseteq \Omega$,
\[
Q(A) - e^\varepsilon P(A) = \sum_{x \in A} \bigl(q(x) - e^\varepsilon p(x)\bigr),
\]
and maximizing termwise gives \eqref{eq:NP2}, attained at $A^* = \{x: q(x) > e^\varepsilon p(x)\}$.
\end{proof}

The two-sided DP curve is
\[
\delta_{\mathrm{two}}(\varepsilon) := \max\{\delta_{Q\|P}(\varepsilon), \delta_{P\|Q}(\varepsilon)\}.
\]

\textbf{Trade-off function.} A complementary view is via the trade-off function (Neyman--Pearson curve)
\[
f_{P,Q}(\alpha) := \inf_{\varphi:\, E_P[\varphi] \leq \alpha} E_Q[1 - \varphi], \quad \alpha \in [0,1],
\]
i.e.\ the minimal type-II error as a function of type-I error.

\subsection{Total variation and Le Cam distance}

For probability measures $\mu, \nu$ on the same measurable space, the total variation distance is
\[
\mathrm{TV}(\mu, \nu) := \sup_A |\mu(A) - \nu(A)| = \tfrac{1}{2}\|\mu - \nu\|_1.
\]

\begin{lemma}[Coupling bound for total variation]
If $(X, Y)$ is any coupling of two random variables with laws $\mu$ and $\nu$, then
\[
\mathrm{TV}(\mu, \nu) \leq P(X \neq Y).
\]
\end{lemma}

\begin{proof}
For any event $A$,
\[
\mu(A) - \nu(A) = P(X \in A) - P(Y \in A) = P(X \in A,\, X \neq Y) - P(Y \in A,\, X \neq Y),
\]
so $|\mu(A) - \nu(A)| \leq P(X \neq Y)$. Take the supremum over $A$.
\end{proof}

\begin{lemma}[Contraction of total variation under measurable maps]
\label{lem:tv-contraction}
Let $\mu, \nu$ be probability measures on $(\mathcal{X}, \mathcal{F})$ and let $f: \mathcal{X} \to \mathcal{Z}$ be measurable. Then
\[
\mathrm{TV}(\mu \circ f^{-1}, \nu \circ f^{-1}) \leq \mathrm{TV}(\mu, \nu).
\]
\end{lemma}

\begin{proof}
For any measurable $B \subseteq \mathcal{Z}$,
\[
(\mu \circ f^{-1})(B) - (\nu \circ f^{-1})(B) = \mu(f^{-1}(B)) - \nu(f^{-1}(B)),
\]
and taking the supremum over $B$ yields a supremum over a subset of measurable sets in $\mathcal{X}$.
\end{proof}

\begin{lemma}[Tensorization / union bound for independent pairs]
\label{lem:tv-tensorization}
Let $(X_1, Y_1)$ and $(X_2, Y_2)$ be couplings such that $X_1 \perp X_2$, $Y_1 \perp Y_2$, and the pair couplings
are independent across indices. Then
\[
\mathrm{TV}(\mathcal{L}(X_1, X_2), \mathcal{L}(Y_1, Y_2)) \leq P(X_1 \neq Y_1) + P(X_2 \neq Y_2).
\]
In particular, if $\mu_i, \nu_i$ are measures on $\mathcal{X}_i$ and $\mu = \mu_1 \otimes \mu_2$, $\nu = \nu_1 \otimes \nu_2$, then
\[
\mathrm{TV}(\mu, \nu) \leq \mathrm{TV}(\mu_1, \nu_1) + \mathrm{TV}(\mu_2, \nu_2).
\]
\end{lemma}

\begin{proof}
Let $(X_1, Y_1)$ and $(X_2, Y_2)$ be couplings realizing the mismatch probabilities. By independence
across indices, $(X_1, X_2)$ is a coupling of $\mu$ and $(Y_1, Y_2)$ is a coupling of $\nu$, and
\[
P\bigl((X_1, X_2) \neq (Y_1, Y_2)\bigr) \leq P(X_1 \neq Y_1) + P(X_2 \neq Y_2)
\]
by the union bound. Apply Lemma~2.2. The product-measure inequality follows by taking couplings
attaining the TV bounds.
\end{proof}

\textbf{Le Cam distance.} Following Le~Cam~\cite{LeC86}, we will use the following elementary bound, sufficient for the present paper.

\begin{lemma}[Le Cam distance bounded by total variation on the same space]
\label{lem:lecam-tv-bound}
Let $(P, Q)$ and $(P', Q')$ be two binary experiments on the same measurable space. Then
\[
\Delta\bigl((P, Q), (P', Q')\bigr) \leq \max\{\mathrm{TV}(P, P'), \mathrm{TV}(Q, Q')\}.
\]
\end{lemma}

\begin{proof}
By definition, the Le Cam distance is the maximum of the two deficiencies. On the same space,
choosing the identity Markov kernel shows each deficiency is bounded by the corresponding total
variation distance, and the claim follows by taking the maximum.
\end{proof}

\begin{lemma}[Privacy curve stability under TV convergence]
\label{lem:privacy-curve-stability}
Let $(P_n, Q_n)$ and $(P, Q)$ be binary experiments. For every $\varepsilon \geq 0$,
\[
\bigl|\delta_{Q_n\|P_n}(\varepsilon) - \delta_{Q\|P}(\varepsilon)\bigr|
\leq \mathrm{TV}(Q_n, Q) + e^\varepsilon\,\mathrm{TV}(P_n, P).
\]
In particular, if $\mathrm{TV}(P_n, P) \to 0$ and $\mathrm{TV}(Q_n, Q) \to 0$, then
$\delta_{Q_n\|P_n}(\varepsilon) \to \delta_{Q\|P}(\varepsilon)$ for every fixed $\varepsilon \geq 0$.
\end{lemma}

\begin{proof}
Using the Neyman--Pearson variational formula \eqref{eq:delta}, for any measurable set $A$,
\begin{align*}
Q_n(A) - e^\varepsilon P_n(A)
&\leq Q(A) + \mathrm{TV}(Q_n,Q) - e^\varepsilon P(A) + e^\varepsilon\mathrm{TV}(P_n,P)\\
&\leq \delta_{Q\|P}(\varepsilon) + \mathrm{TV}(Q_n,Q) + e^\varepsilon\mathrm{TV}(P_n,P).
\end{align*}
Taking the supremum over $A$ gives $\delta_{Q_n\|P_n}(\varepsilon) \leq \delta_{Q\|P}(\varepsilon) + \mathrm{TV}(Q_n,Q) + e^\varepsilon\mathrm{TV}(P_n,P)$.
The reverse inequality is symmetric.
\end{proof}

\subsection{Binary randomized response and the scaling parameter}

In Sections~3--4 we specialize to binary randomized response (RR) with an $n$-dependent local privacy
level $\varepsilon_0 = \varepsilon_0(n) > 0$. The local channel $W^{(n)}$ is
\[
W^{(n)}_0(1) = \delta_n, \quad W^{(n)}_1(1) = 1 - \delta_n, \quad \delta_n := \frac{1}{1 + e^{\varepsilon_0(n)}}.
\]
Equivalently, the local likelihood ratios are $w(1) = e^{\varepsilon_0(n)}$ and $w(0) = e^{-\varepsilon_0(n)}$. A key scaling
parameter is
\begin{equation}
a_n := \frac{e^{\varepsilon_0(n)}}{n}. \label{eq:an}
\end{equation}
The sub-critical / critical / super-critical regimes correspond to $a_n \to 0$, $a_n \to c^2 \in (0, \infty)$, and
$a_n \to \infty$, respectively.

\section{Canonical neighboring pair: Poisson-shift limit}

We first treat the canonical neighboring datasets: all zeros versus one one, under shuffled RR with
$n$-dependent $\varepsilon_0(n)$.

\subsection{Exact one-dimensional reduction}

Let $K_n$ be the released count of output-1 messages. Under the canonical null (all zeros),
\begin{equation}
K_n \sim \mathrm{Bin}(n, \delta_n). \label{eq:Kn_null}
\end{equation}
Under the canonical alternative (one one),
\begin{equation}
K_n \sim \mathrm{Bin}(n-1, \delta_n) + \mathrm{Bern}(1 - \delta_n), \label{eq:Kn_alt}
\end{equation}
with independence between the two terms.

The shuffle likelihood ratio depends on the transcript only through $K_n$ and admits the explicit
affine form
\begin{equation}
L_n(K_n) = \frac{1}{n}\Bigl[(n - K_n)e^{-\varepsilon_0(n)} + K_n e^{\varepsilon_0(n)}\Bigr]
= e^{-\varepsilon_0(n)} + \frac{e^{\varepsilon_0(n)} - e^{-\varepsilon_0(n)}}{n} K_n. \label{eq:LR}
\end{equation}

\textit{Derivation of \eqref{eq:LR}.} Condition on the histogram (equivalently on $K_n$). Under the null, all $n$ users
draw i.i.d.\ from $W^{(n)}_0$; under the alternative, one user draws from $W^{(n)}_1$ and the remaining $n-1$
users draw from $W^{(n)}_0$. Given the multiset, the identity of the ``special'' message is uniformly random
among the $n$ messages. Therefore the likelihood ratio is the average, over the $n$ positions, of the
local likelihood ratio $W^{(n)}_1(y)/W^{(n)}_0(y)$ evaluated at the message in that position. For RR, this ratio
equals $e^{\varepsilon_0(n)}$ on output 1 and $e^{-\varepsilon_0(n)}$ on output 0. Among the $n$ messages, $K_n$ are equal to 1 and
$n - K_n$ are equal to 0, yielding \eqref{eq:LR}.

\subsection{Critical scaling and Poisson approximation}

Assume the critical regime
\begin{equation}
a_n = \frac{e^{\varepsilon_0(n)}}{n} \to c^2 \in (0, \infty). \label{eq:critical}
\end{equation}
Then $\delta_n = (1 + e^{\varepsilon_0(n)})^{-1} \asymp n^{-1}$ and the binomial counts in \eqref{eq:Kn_null}--\eqref{eq:Kn_alt}
have $O(1)$ means. The correct limit is Poisson.

We will use two explicit approximation lemmas, proved in Appendix~A: a binomial-to-Poisson
total-variation bound (Lemma~A.1) and a Poisson parameter perturbation bound (Lemma~A.2).

\begin{theorem}[Poisson-shift limit experiment]
\label{thm:poisson-shift}
Assume the critical regime \eqref{eq:critical}. Let $P_n$ be the law of $K_n$ under \eqref{eq:Kn_null} and $Q_n$ be the
law of $K_n$ under \eqref{eq:Kn_alt}. Let $\lambda := c^{-2}$ and define the Poisson-shift limit experiment
\[
\mathcal{E}^{\mathrm{Poi}}_\infty := (P_\infty, Q_\infty) := \bigl(\mathrm{Poi}(\lambda),\; 1 + \mathrm{Poi}(\lambda)\bigr).
\]
Then $\Delta\bigl((P_n, Q_n), (P_\infty, Q_\infty)\bigr) \to 0$.

Moreover, letting $\lambda_n := n\delta_n$ and $\lambda_{n-1} := (n-1)\delta_n$, we have the explicit bounds
\begin{align}
\mathrm{TV}(P_n, P_\infty) &\leq n\delta_n(1 - e^{-\delta_n}) + |\lambda_n - \lambda| \leq n\delta_n^2 + |\lambda_n - \lambda|, \label{eq:TVPn}\\
\mathrm{TV}(Q_n, Q_\infty) &\leq \delta_n + (n-1)\delta_n(1 - e^{-\delta_n}) + |\lambda_{n-1} - \lambda| \leq \delta_n + (n-1)\delta_n^2 + |\lambda_{n-1} - \lambda|. \label{eq:TVQn}
\end{align}
Consequently,
\begin{equation}
\Delta\bigl((P_n, Q_n), (P_\infty, Q_\infty)\bigr) \leq \max\{\mathrm{TV}(P_n, P_\infty), \mathrm{TV}(Q_n, Q_\infty)\}. \label{eq:LeCam_bound}
\end{equation}
In particular, in the canonical calibration $e^{\varepsilon_0(n)} = c^2 n$ (equivalently $\delta_n = (1 + c^2 n)^{-1}$),
\begin{equation}
\Delta\bigl((P_n, Q_n), (P_\infty, Q_\infty)\bigr) \leq \frac{2}{c^2 n} + \frac{2}{c^4 n}. \label{eq:explicit_rate}
\end{equation}
\end{theorem}

\begin{proof}
\textbf{Step 1: approximate the null $\mathrm{Bin}(n, \delta_n)$ by $\mathrm{Poi}(\lambda_n)$.} Let $S_n \sim \mathrm{Bin}(n, \delta_n)$ and let
$N_n \sim \mathrm{Poi}(\lambda_n)$ with $\lambda_n = n\delta_n$. By Lemma~A.1,
\[
\mathrm{TV}(\mathcal{L}(S_n), \mathrm{Poi}(\lambda_n)) \leq n\delta_n(1 - e^{-\delta_n}) \leq n\delta_n^2.
\]

\textbf{Step 2: perturb $\mathrm{Poi}(\lambda_n)$ to $\mathrm{Poi}(\lambda)$.} By Lemma~A.2,
\[
\mathrm{TV}(\mathrm{Poi}(\lambda_n), \mathrm{Poi}(\lambda)) \leq 1 - e^{-|\lambda_n - \lambda|} \leq |\lambda_n - \lambda|.
\]
By the triangle inequality this yields \eqref{eq:TVPn}.

\textbf{Step 3: approximate the alternative.} Under $Q_n$, we may write $K_n = S_{n-1} + B_n$ with
$S_{n-1} \sim \mathrm{Bin}(n-1, \delta_n)$ and $B_n \sim \mathrm{Bern}(1 - \delta_n)$ independent. Couple $B_n$ to the constant 1 so that
$P(B_n \neq 1) = \delta_n$. By Lemma~2.2 and the shift $K_n^{\mathrm{alt}} = S_{n-1} + B_n$,
\[
\mathrm{TV}\bigl(\mathcal{L}(K_n^{\mathrm{alt}}), \mathcal{L}(S_{n-1} + 1)\bigr) \leq \delta_n.
\]

Next, apply Lemma~A.1 with $m = n-1$ to couple $S_{n-1}$ to $\mathrm{Poi}(\lambda_{n-1})$:
\[
\mathrm{TV}(\mathcal{L}(S_{n-1}), \mathrm{Poi}(\lambda_{n-1})) \leq (n-1)\delta_n(1 - e^{-\delta_n}) \leq (n-1)\delta_n^2.
\]
Finally, apply Lemma~A.2 to perturb $\mathrm{Poi}(\lambda_{n-1})$ to $\mathrm{Poi}(\lambda)$:
\[
\mathrm{TV}(\mathrm{Poi}(\lambda_{n-1}), \mathrm{Poi}(\lambda)) \leq |\lambda_{n-1} - \lambda|.
\]
Combine these three steps by the triangle inequality (and shift invariance) to obtain \eqref{eq:TVQn}.

\textbf{Step 4: Le Cam distance.} Since both experiments live on the same countable space $\mathbb{N}$,
Lemma~2.5 gives \eqref{eq:LeCam_bound}. Under \eqref{eq:critical}, we have $\delta_n \to 0$ and $\lambda_n, \lambda_{n-1} \to \lambda$, so the right-hand side
converges to 0.

\textbf{Step 5: explicit $O(n^{-1})$ constant under canonical calibration.} Assume $e^{\varepsilon_0(n)} = c^2 n$, so
$\delta_n = (1 + c^2 n)^{-1}$. Then
\[
n\delta_n^2 = \frac{n}{(1 + c^2 n)^2} \leq \frac{1}{c^4 n}, \qquad
|\lambda_n - \lambda| = \left|\frac{n}{1 + c^2 n} - \frac{1}{c^2}\right| = \frac{1}{c^2(1 + c^2 n)} \leq \frac{1}{c^4 n},
\]
hence $\mathrm{TV}(P_n, P_\infty) \leq 2/(c^4 n)$.

Similarly, $\delta_n \leq 1/(c^2 n)$, $(n-1)\delta_n^2 \leq 1/(c^4 n)$, and
\[
|\lambda_{n-1} - \lambda| = \left|\frac{n-1}{1 + c^2 n} - \frac{1}{c^2}\right| = \frac{1 + c^2}{c^2(1 + c^2 n)} \leq \frac{1 + c^2}{c^4 n}.
\]
Therefore
\[
\mathrm{TV}(Q_n, Q_\infty) \leq \frac{1}{c^2 n} + \frac{1}{c^4 n} + \frac{1 + c^2}{c^4 n} = \frac{2}{c^2 n} + \frac{2}{c^4 n}.
\]
Since $\mathrm{TV}(P_n, P_\infty) \leq 2/(c^4 n) \leq 2/(c^2 n) + 2/(c^4 n) = \mathrm{TV}(Q_n, Q_\infty)$, both are bounded by $2/(c^2 n) + 2/(c^4 n)$, and \eqref{eq:explicit_rate} follows from \eqref{eq:LeCam_bound}.
\end{proof}

\begin{corollary}[Privacy curve convergence for the Poisson-shift limit]
\label{cor:poisson-privacy-curve-convergence}
Under the assumptions of Theorem~\ref{thm:poisson-shift}, for every fixed $\varepsilon \geq 0$,
\[
\delta_{Q_n\|P_n}(\varepsilon) \to \delta_{Q_\infty\|P_\infty}(\varepsilon).
\]
Under the canonical calibration $e^{\varepsilon_0(n)} = c^2 n$, the convergence is quantitative:
\[
\bigl|\delta_{Q_n\|P_n}(\varepsilon) - \delta_{Q_\infty\|P_\infty}(\varepsilon)\bigr|
\leq (1 + e^\varepsilon)\Bigl(\frac{2}{c^2 n} + \frac{2}{c^4 n}\Bigr).
\]
\end{corollary}

\begin{proof}
Apply Lemma~\ref{lem:privacy-curve-stability} with the TV bounds \eqref{eq:TVPn}--\eqref{eq:TVQn}
and the explicit rate \eqref{eq:explicit_rate}.
\end{proof}

\begin{proposition}[Sharp $n^{-1}$ total-variation rate in the canonical Poisson regime]
\label{prop:poisson-tv-rate}
Assume the setting of Theorem~\ref{thm:poisson-shift} and the canonical calibration $e^{\varepsilon_0(n)}=c^2 n$, so that
$\delta_n=(1+c^2 n)^{-1}$ and $\lambda=c^{-2}$. Then, for all sufficiently large $n$,
\[
\frac{e^{-1/c^2}}{4c^4}\,\frac1n
\le \mathrm{TV}(P_n,P_\infty)
\le \frac{2}{c^4}\,\frac1n.
\]
In particular,
\[
\mathrm{TV}(P_n,P_\infty)=\Theta(n^{-1}).
\]
More precisely,
\[
P_n\{0\}-P_\infty\{0\}
= \frac{e^{-1/c^2}}{2c^4}\,\frac1n + O(n^{-2})
= \frac{e^{-1/c^2}}{2}\,n\delta_n^2 + O(n^{-2}).
\]
\end{proposition}

\begin{proof}
The upper bound is exactly the $P_n$-bound from Theorem~\ref{thm:poisson-shift} under the canonical calibration.
For the lower bound, let $p_n(k):=P_n\{k\}$ and $p_\infty(k):=P_\infty\{k\}$. Then
\[
p_n(0)=(1-\delta_n)^n=\left(1+\frac{1}{c^2 n}\right)^{-n},
\qquad
p_\infty(0)=e^{-1/c^2}.
\]
Hence
\[
\log\frac{p_n(0)}{p_\infty(0)}
=\frac1{c^2}-n\log\left(1+\frac1{c^2 n}\right)
=\frac{1}{2c^4 n}-\frac{1}{3c^6 n^2}+O(n^{-3}),
\]
so
\[
p_n(0)-p_\infty(0)
=e^{-1/c^2}\Bigg(\exp\!\left(\frac{1}{2c^4 n}+O(n^{-2})\right)-1\Bigg)
=\frac{e^{-1/c^2}}{2c^4}\,\frac1n+O(n^{-2}).
\]
Since $\mathrm{TV}(P_n,P_\infty)\ge |p_n(0)-p_\infty(0)|$, this yields the lower bound and hence the claim.
\end{proof}

\begin{remark}[Chen--Stein comparison]
For sums of independent Bernoulli indicators, the classical Chen--Stein / Le Cam bound gives
\[
\mathrm{TV}\!\left(\mathcal{L}\!\left(\sum_i I_i\right), \mathrm{Poi}(\lambda)\right) \leq \sum_i p_i^2, \quad \lambda := \sum_i p_i.
\]
For $S \sim \mathrm{Bin}(m, p)$ this becomes $\mathrm{TV}(\mathcal{L}(S), \mathrm{Poi}(mp)) \leq mp^2$, which matches the order delivered by
Lemma~A.1 and by \eqref{eq:TVPn}--\eqref{eq:TVQn}. Thus Stein's method does not improve the leading $O(mp^2)$ rate in
the present i.i.d.\ setting, but it does provide a flexible route to non-identically distributed or weakly
dependent rare-event arrays; see \cite{BHJ92, CGS11}. For translated Poisson refinements, see also \cite{Rol07}.
\end{remark}

\subsection{Limiting privacy curve and the \texorpdfstring{$\delta$}{delta}-floor}

\begin{proposition}[Limiting curve as a Poisson series]
\label{prop:poisson-curve}
Let $P_\infty = \mathrm{Poi}(\lambda)$ and $Q_\infty$ be the law of $1 + J$ with $J \sim \mathrm{Poi}(\lambda)$, where $\lambda = c^{-2}$.
Then for every $\varepsilon \geq 0$,
\begin{equation}
\delta_{Q_\infty\|P_\infty}(\varepsilon) = \sum_{j \geq 0} \bigl(P_\infty(j-1) - e^\varepsilon P_\infty(j)\bigr)_+, \quad P_\infty(-1) := 0. \label{eq:Poi_curve}
\end{equation}
Moreover, the two-sided curve has a support-mismatch floor:
\begin{equation}
\delta_{\mathrm{two}}(\varepsilon) \geq \delta_{P_\infty\|Q_\infty}(\varepsilon) \geq P_\infty(\{0\}) = e^{-\lambda}, \qquad \forall\, \varepsilon \geq 0. \label{eq:floor}
\end{equation}
\end{proposition}

\begin{proof}
Since $P_\infty$ and $Q_\infty$ are supported on $\mathbb{N}$ and $Q_\infty(j) = P_\infty(j-1)$ for $j \geq 1$ (and $Q_\infty(0) = 0$),
the discrete identity in Lemma~2.1 yields \eqref{eq:Poi_curve}.

For the floor, note that $Q_\infty(\{0\}) = 0$ while $P_\infty(\{0\}) = e^{-\lambda}$. Taking $A = \{0\}$ in the definition \eqref{eq:delta}
of $\delta_{P_\infty\|Q_\infty}(\varepsilon)$ gives
\[
\delta_{P_\infty\|Q_\infty}(\varepsilon) \geq P_\infty(\{0\}) - e^\varepsilon Q_\infty(\{0\}) = e^{-\lambda}.
\]
\end{proof}

\begin{remark}[Non-commuting limits: the Poisson floor is a large-$n$ phenomenon]
\label{rem:noncommuting-limits}
For each finite $n$, both $P_n = \mathrm{Bin}(n,\delta_n)$ and $Q_n$ have full support on $\mathbb{N}$,
so the likelihood ratio $L_n$ is strictly positive everywhere and
\[
\lim_{\varepsilon \to \infty} \delta^{(n)}_{\mathrm{two}}(\varepsilon) = 0 \qquad \text{for every fixed } n.
\]
The floor $e^{-\lambda}$ in \eqref{eq:floor} belongs to the \emph{limit experiment} and arises because
$Q_\infty(\{0\}) = 0$: the atom at zero has support-mismatch in the limit. Therefore the two iterated
limits do not commute:
\[
\lim_{\varepsilon \to \infty}\lim_{n \to \infty} \delta^{(n)}_{\mathrm{two}}(\varepsilon)
= e^{-\lambda} > 0,
\qquad
\lim_{n \to \infty}\lim_{\varepsilon \to \infty} \delta^{(n)}_{\mathrm{two}}(\varepsilon) = 0.
\]
This non-commutativity is a genuine feature of the critical regime and not a defect of the approximation:
the floor $e^{-\lambda}$ is an intrinsic property of the Poisson-shift limit experiment capturing
the probability of zero errors, an event that separates the two hypotheses with probability $e^{-\lambda}$
regardless of the threshold $e^\varepsilon$.
\end{remark}

\begin{proposition}[Trade-off function of the Poisson-shift limit and recovery of the privacy curve]
\label{prop:poisson-tradeoff}
Let $(P_\infty,Q_\infty)$ be the Poisson-shift limit experiment of Theorem~\ref{thm:poisson-shift}, so that
\[
P_\infty=\mathrm{Poi}(\lambda),
\qquad
Q_\infty=1+\mathrm{Poi}(\lambda),
\qquad \lambda=c^{-2}>0,
\]
and let $J\sim \mathrm{Poi}(\lambda)$. Write
\[
p_\lambda(m):=\mathbb P(J=m)=e^{-\lambda}\frac{\lambda^m}{m!},
\qquad
\overline F_\lambda(m):=\mathbb P(J\ge m),\quad m\in\mathbb N.
\]
Then the trade-off function $f_{P_\infty,Q_\infty}$ is piecewise affine with knots at the Poisson upper-tail levels $\overline F_\lambda(m)$:
\[
f_{P_\infty,Q_\infty}(0)=1,
\qquad
f_{P_\infty,Q_\infty}(\alpha)=0\quad\text{for }\alpha\in[\overline F_\lambda(1),1],
\]
and, for every $m\ge 1$ and every
\[
\overline F_\lambda(m+1)\le \alpha\le \overline F_\lambda(m),
\]
we have
\[
f_{P_\infty,Q_\infty}(\alpha)
=1-\overline F_\lambda(m)-\frac{m}{\lambda}\bigl(\alpha-\overline F_\lambda(m+1)\bigr)
=\mathbb P(J\le m-1)-\frac{m}{\lambda}\bigl(\alpha-\mathbb P(J\ge m+1)\bigr).
\]
Equivalently, the Neyman--Pearson curve is the lower convex envelope of the points
\[
\bigl(\mathbb P(J\ge m),\, \mathbb P(J\le m-2)\bigr),\qquad m\ge 1,
\]
with the convention $\mathbb P(J\le -1)=0$.

Moreover, for every $\varepsilon\ge 0$,
\[
\delta_{Q_\infty\|P_\infty}(\varepsilon)
=\sup_{\alpha\in[0,1]}\Bigl\{1-f_{P_\infty,Q_\infty}(\alpha)-e^\varepsilon\alpha\Bigr\}.
\]
If $m_\varepsilon:=\lfloor \lambda e^\varepsilon\rfloor+1$, then one maximizer is
\[
\alpha_\varepsilon=\mathbb P(J\ge m_\varepsilon),
\]
and therefore
\[
\delta_{Q_\infty\|P_\infty}(\varepsilon)
=\mathbb P(J\ge m_\varepsilon-1)-e^\varepsilon\mathbb P(J\ge m_\varepsilon).
\]
This is exactly the Poisson upper-tail form of Proposition~\ref{prop:poisson-curve}.
\end{proposition}

\begin{proof}
Let
\[
p(j):=P_\infty\{j\}=e^{-\lambda}\frac{\lambda^j}{j!},
\qquad
q(j):=Q_\infty\{j\}=p(j-1)\quad (j\ge 1),
\qquad q(0)=0.
\]
Then the likelihood ratio is
\[
L(j):=\frac{q(j)}{p(j)}=
\begin{cases}
0, & j=0,\\[2mm]
\dfrac{j}{\lambda}, & j\ge 1.
\end{cases}
\]
Hence $L(j)$ is strictly increasing in $j$. By the Neyman--Pearson lemma, the most powerful tests for
$P_\infty$ versus $Q_\infty$ are upper-tail tests with possible randomization at one boundary point:
for some $m\ge 1$ and $\tau\in[0,1]$,
\[
\varphi_{m,\tau}(j)=\mathbf 1\{j\ge m+1\}+\tau\,\mathbf 1\{j=m\}.
\]
Its type-I error is
\[
\alpha=\mathbb E_{P_\infty}[\varphi_{m,\tau}]
=\mathbb P(J\ge m+1)+\tau\,\mathbb P(J=m)
=\overline F_\lambda(m+1)+\tau p_\lambda(m),
\]
so necessarily $\alpha\in[\overline F_\lambda(m+1),\overline F_\lambda(m)]$, and
\[
\tau=\frac{\alpha-\overline F_\lambda(m+1)}{p_\lambda(m)}.
\]
The corresponding type-II error is
\begin{align*}
f_{P_\infty,Q_\infty}(\alpha)
&=\mathbb E_{Q_\infty}[1-\varphi_{m,\tau}]\\
&=1-Q_\infty\{j\ge m+1\}-\tau Q_\infty\{j=m\}\\
&=1-\mathbb P(J\ge m)-\tau\,\mathbb P(J=m-1).
\end{align*}
Since
\[
\frac{\mathbb P(J=m-1)}{\mathbb P(J=m)}=\frac{m}{\lambda},
\]
eliminating $\tau$ gives
\[
f_{P_\infty,Q_\infty}(\alpha)
=1-\overline F_\lambda(m)-\frac{m}{\lambda}\bigl(\alpha-\overline F_\lambda(m+1)\bigr),
\]
which proves the stated piecewise-affine formula. If $\alpha\ge \overline F_\lambda(1)=\mathbb P(J\ge 1)$,
we may reject on $\{j\ge 1\}$ and randomize further at $j=0$; since $Q_\infty\{0\}=0$, the resulting
minimal type-II error is $0$. Also $f_{P_\infty,Q_\infty}(0)=1$.

For the relation with the privacy curve, let $\varphi$ be any test, with type-I error
$\alpha(\varphi):=\mathbb E_{P_\infty}[\varphi]$ and type-II error
$\beta(\varphi):=\mathbb E_{Q_\infty}[1-\varphi]$. Then
\[
Q_\infty(\varphi=1)-e^\varepsilon P_\infty(\varphi=1)
=1-\beta(\varphi)-e^\varepsilon\alpha(\varphi).
\]
Allowing randomized tests does not change the supremum in the privacy-curve variational formula, since the objective is linear in $\varphi$. Taking the supremum over all tests and then minimizing over tests at fixed level $\alpha$ yields the standard identity
\[
\delta_{Q_\infty\|P_\infty}(\varepsilon)
=\sup_{\alpha\in[0,1]}\Bigl\{1-f_{P_\infty,Q_\infty}(\alpha)-e^\varepsilon\alpha\Bigr\}.
\]
Now fix $\varepsilon\ge 0$. On the interval
$[\overline F_\lambda(m+1),\overline F_\lambda(m)]$, the function
\[
\alpha\longmapsto 1-f_{P_\infty,Q_\infty}(\alpha)-e^\varepsilon\alpha
\]
is affine with slope
\[
\frac{m}{\lambda}-e^\varepsilon.
\]
Hence it is increasing when $m/\lambda>e^\varepsilon$ and decreasing when $m/\lambda<e^\varepsilon$. If $\lambda e^\varepsilon$ is an integer, the maximum is attained on an entire boundary segment; the choice below is one convenient maximizer. Thus a maximizer is attained at the knot corresponding to
\[
m_\varepsilon=\lfloor \lambda e^\varepsilon\rfloor+1,
\qquad
\alpha_\varepsilon=\overline F_\lambda(m_\varepsilon)=\mathbb P(J\ge m_\varepsilon).
\]
Evaluating at that point gives
\begin{align*}
\delta_{Q_\infty\|P_\infty}(\varepsilon)
&=1-f_{P_\infty,Q_\infty}(\alpha_\varepsilon)-e^\varepsilon\alpha_\varepsilon\\
&=1-\mathbb P(J\le m_\varepsilon-2)-e^\varepsilon\mathbb P(J\ge m_\varepsilon)\\
&=\mathbb P(J\ge m_\varepsilon-1)-e^\varepsilon\mathbb P(J\ge m_\varepsilon).
\end{align*}
Since $q(j)-e^\varepsilon p(j)=p(j)\bigl(j/\lambda-e^\varepsilon\bigr)$, the active set in the discrete
Neyman--Pearson formula of Lemma~2.1 is exactly $\{j\ge m_\varepsilon\}$; thus the last display is the
same formula as Proposition~\ref{prop:poisson-curve}, written as a Poisson upper-tail difference.
\end{proof}

\begin{proposition}[Monotonicity of the Poisson-shift limiting curve]
\label{prop:poisson-monotonicity}
For $\lambda > 0$, let $P_\lambda := \mathrm{Poi}(\lambda)$ and $Q_\lambda := 1 + \mathrm{Poi}(\lambda)$.
\begin{enumerate}
\item[(i)] For fixed $\lambda$, the one-sided curve $\varepsilon \mapsto \delta_{Q_\lambda\|P_\lambda}(\varepsilon)$ is strictly decreasing on $[0, \infty)$.
\item[(ii)] For fixed $\varepsilon \geq 0$, the map $\lambda \mapsto \delta_{Q_\lambda\|P_\lambda}(\varepsilon)$ is strictly decreasing on $(0, \infty)$. Equivalently, under
the paper's parameterization $\lambda = c^{-2}$, the map $c \mapsto \delta_{Q_c\|P_c}(\varepsilon)$ is strictly increasing: larger $c$
means weaker privacy.
\item[(iii)] For fixed $\lambda$, the reverse curve $\varepsilon \mapsto \delta_{P_\lambda\|Q_\lambda}(\varepsilon)$ is strictly decreasing until it reaches the floor $e^{-\lambda}$,
after which it is constant. Hence the two-sided Poisson limiting curve is nonincreasing in $\varepsilon$
and bounded below by $e^{-\lambda}$.
\end{enumerate}
\end{proposition}

\begin{proof}
For (i), write $L(j) := Q_\lambda(j)/P_\lambda(j) = j/\lambda$ for $j \geq 1$ and $L(0) = 0$. Since $L(j) \to \infty$ as
$j \to \infty$, for every $\varepsilon_2 > \varepsilon_1 \geq 0$ the pointwise inequality $(L - e^{\varepsilon_2})_+ \leq (L - e^{\varepsilon_1})_+$ is strict on the event
$\{L > e^{\varepsilon_2}\}$, which has positive $P_\lambda$-probability. Taking expectations under $P_\lambda$ proves strict decrease.

For (ii), let $m(\lambda, \varepsilon) := \lfloor \lambda e^\varepsilon \rfloor + 1$. The Neyman--Pearson set is the upper tail $\{j \geq m(\lambda, \varepsilon)\}$, so
\[
\delta_{Q_\lambda\|P_\lambda}(\varepsilon) = P_\lambda(J \geq m - 1) - e^\varepsilon P_\lambda(J \geq m), \quad J \sim \mathrm{Poi}(\lambda).
\]
On each interval of $\lambda$ for which $m$ is constant, differentiating the Poisson tails gives
\[
\frac{d}{d\lambda}\delta_{Q_\lambda\|P_\lambda}(\varepsilon) = \begin{cases} -e^\varepsilon P_\lambda(J = 0), & m = 1, \\ P_\lambda(J = m-2) - e^\varepsilon P_\lambda(J = m-1), & m \geq 2. \end{cases}
\]
If $m \geq 2$, then $m - 1 \leq \lambda e^\varepsilon$, hence
\[
P_\lambda(J = m-2) - e^\varepsilon P_\lambda(J = m-1) = P_\lambda(J = m-2)\!\left(1 - \frac{e^\varepsilon \lambda}{m-1}\right) < 0
\]
away from the breakpoints. At a breakpoint $\lambda_0$, the series formula of Proposition~\ref{prop:poisson-curve} shows that $\delta_{Q_\lambda\|P_\lambda}(\varepsilon)$ is
continuous in $\lambda$, by dominated convergence applied termwise to the Poisson pmf series. Hence strict decrease on each
open interval on which $m$ is constant, together with continuity at the breakpoints, yields global strict decrease in $\lambda$.

For (iii), note that $P_\lambda(j)/Q_\lambda(j) = \lambda/j$ for $j \geq 1$, so the active set in the discrete Neyman--Pearson
formula is finite. Once $e^\varepsilon \geq \lambda$, no $j \geq 1$ contributes and only the atom at 0 remains, giving the
constant floor $e^{-\lambda}$. Before that threshold, the same strict-decrease argument as in (i) applies.
\end{proof}

\section{Proportional compositions: Skellam-shift limit}

We now consider compositions $k = k(n)$ with $\pi_n := k/n \to \pi \in (0, 1)$ under shuffled RR in the
critical scaling $a_n \to c^2$.

\subsection{Centered decomposition}

Let $K_{n,k}$ be the released number of ones under composition $k$. Let
\begin{align*}
A_{n,k} &\sim \mathrm{Bin}(n-k, \delta_n) \quad \text{(false positives from 0-users)},\\
B_{n,k} &\sim \mathrm{Bin}(k, \delta_n) \quad \text{(false negatives from 1-users)},
\end{align*}
independent. Then
\begin{equation}
K_{n,k} = k + A_{n,k} - B_{n,k}, \qquad D_{n,k} := K_{n,k} - k = A_{n,k} - B_{n,k}. \label{eq:Dnk}
\end{equation}
The neighboring dataset with $k+1$ ones satisfies the analogous decomposition
\begin{equation}
K_{n,k+1} = k + 1 + A'_{n,k} - B'_{n,k}, \qquad D'_{n,k} := K_{n,k+1} - k = 1 + A'_{n,k} - B'_{n,k}, \label{eq:Dnk_alt}
\end{equation}
where $A'_{n,k} \sim \mathrm{Bin}(n-k-1, \delta_n)$ and $B'_{n,k} \sim \mathrm{Bin}(k+1, \delta_n)$ are independent.

\subsection{Skellam limit}

\begin{theorem}[Skellam-shift limit experiment]
\label{thm:skellam-shift}
Assume $\pi_n = k/n \to \pi \in (0,1)$ and $a_n = e^{\varepsilon_0(n)}/n \to c^2 \in (0, \infty)$. Define
\[
\lambda_0 := \frac{1 - \pi}{c^2}, \qquad \lambda_1 := \frac{\pi}{c^2}.
\]
Let $P_n$ be the law of $D_{n,k}$ in \eqref{eq:Dnk} under $T_{n,k}$ and $Q_n$ be the law of $D'_{n,k}$ in \eqref{eq:Dnk_alt} under $T_{n,k+1}$. Let
$D \sim \mathrm{Skellam}(\lambda_0, \lambda_1)$, i.e.\ $D = X - Y$ with $X \sim \mathrm{Poi}(\lambda_0)$ and $Y \sim \mathrm{Poi}(\lambda_1)$ independent, and define
the Skellam-shift limit experiment
\[
\mathcal{E}^{\mathrm{Skellam}}_\infty := (P_\infty, Q_\infty) := \bigl(\mathcal{L}(D), \mathcal{L}(1 + D)\bigr).
\]
Then $P_n \to P_\infty$ and $Q_n \to Q_\infty$ in total variation, hence $\Delta\bigl((P_n, Q_n), (P_\infty, Q_\infty)\bigr) \to 0$.

Moreover, letting $\lambda_{0,n} := (n-k)\delta_n$ and $\lambda_{1,n} := k\delta_n$, we have the explicit bounds
\begin{align}
\mathrm{TV}(P_n, P_\infty) &\leq (n-k)\delta_n(1-e^{-\delta_n}) + k\delta_n(1-e^{-\delta_n}) + |\lambda_{0,n} - \lambda_0| + |\lambda_{1,n} - \lambda_1| \nonumber\\
&\leq n\delta_n^2 + |\lambda_{0,n} - \lambda_0| + |\lambda_{1,n} - \lambda_1|, \label{eq:TVPn_Sk}\\
\mathrm{TV}(Q_n, Q_\infty) &\leq (n-k-1)\delta_n(1-e^{-\delta_n}) + (k+1)\delta_n(1-e^{-\delta_n}) + |\lambda'_{0,n} - \lambda_0| + |\lambda'_{1,n} - \lambda_1| \nonumber\\
&\leq n\delta_n^2 + |\lambda'_{0,n} - \lambda_0| + |\lambda'_{1,n} - \lambda_1|, \label{eq:TVQn_Sk}
\end{align}
where $\lambda'_{0,n} := (n-k-1)\delta_n$ and $\lambda'_{1,n} := (k+1)\delta_n$.
\end{theorem}

\begin{proof}
\textbf{Step 1: Poisson approximation for $A_{n,k}$ and $B_{n,k}$.} By Lemma~A.1 applied to $A_{n,k} \sim \mathrm{Bin}(n-k, \delta_n)$, we can couple $A_{n,k}$ to $\tilde{A}_{n,k} \sim \mathrm{Poi}(\lambda_{0,n})$ so that
\[
\mathrm{TV}(\mathcal{L}(A_{n,k}), \mathrm{Poi}(\lambda_{0,n})) \leq (n-k)\delta_n(1-e^{-\delta_n}).
\]
Similarly,
\[
\mathrm{TV}(\mathcal{L}(B_{n,k}), \mathrm{Poi}(\lambda_{1,n})) \leq k\delta_n(1-e^{-\delta_n}).
\]

\textbf{Step 2: push the approximation through the difference map.} Let $\tilde{B}_{n,k} \sim \mathrm{Poi}(\lambda_{1,n})$
independent of $\tilde{A}_{n,k}$. By Lemma~2.4,
\[
\mathrm{TV}\bigl(\mathcal{L}(A_{n,k}, B_{n,k}), \mathcal{L}(\tilde{A}_{n,k}, \tilde{B}_{n,k})\bigr) \leq (n-k)\delta_n(1-e^{-\delta_n}) + k\delta_n(1-e^{-\delta_n}).
\]
Since total variation cannot increase under the measurable map $(a,b) \mapsto a - b$, we obtain
\begin{equation}
\mathrm{TV}\bigl(\mathcal{L}(D_{n,k}), \mathcal{L}(\tilde{A}_{n,k} - \tilde{B}_{n,k})\bigr) \leq (n-k)\delta_n(1-e^{-\delta_n}) + k\delta_n(1-e^{-\delta_n}). \label{eq:step2}
\end{equation}

\textbf{Step 3: perturb Poisson parameters to the limits.} Let $X \sim \mathrm{Poi}(\lambda_0)$ and $Y \sim \mathrm{Poi}(\lambda_1)$
independent, so $D = X - Y \sim \mathrm{Skellam}(\lambda_0, \lambda_1)$. By Lemma~A.2,
\[
\mathrm{TV}(\mathrm{Poi}(\lambda_{0,n}), \mathrm{Poi}(\lambda_0)) \leq |\lambda_{0,n} - \lambda_0|, \qquad \mathrm{TV}(\mathrm{Poi}(\lambda_{1,n}), \mathrm{Poi}(\lambda_1)) \leq |\lambda_{1,n} - \lambda_1|.
\]
Using Lemma~2.4 and contraction under $(x,y) \mapsto x - y$ again gives
\begin{equation}
\mathrm{TV}\bigl(\mathcal{L}(\tilde{A}_{n,k} - \tilde{B}_{n,k}), \mathcal{L}(X - Y)\bigr) \leq |\lambda_{0,n} - \lambda_0| + |\lambda_{1,n} - \lambda_1|. \label{eq:step3}
\end{equation}
Combining \eqref{eq:step2} and \eqref{eq:step3} by the triangle inequality yields \eqref{eq:TVPn_Sk}. The final inequality uses $1 - e^{-\delta_n} \leq \delta_n$.

\textbf{Step 4: the alternative.} The same argument applies to $A'_{n,k} \sim \mathrm{Bin}(n-k-1, \delta_n)$ and
$B'_{n,k} \sim \mathrm{Bin}(k+1, \delta_n)$, yielding \eqref{eq:TVQn_Sk} for the law of $A'_{n,k} - B'_{n,k}$. Since $D'_{n,k} = 1 + (A'_{n,k} - B'_{n,k})$ and
shifts preserve total variation, the same bound holds for $Q_n$ relative to $Q_\infty = \mathcal{L}(1 + D)$.

\textbf{Step 5: Le Cam distance and convergence.} Both experiments live on the countable space $\mathbb{Z}$,
so Lemma~2.5 yields experiment-level convergence. Under $a_n \to c^2$ and $\pi_n \to \pi$, we have $\lambda_{0,n} \to \lambda_0$
and $\lambda_{1,n} \to \lambda_1$, while $n\delta_n^2 \to 0$; thus the TV bounds vanish.
\end{proof}

\begin{corollary}[Explicit $O(n^{-1})$ rate under canonical calibration]
\label{cor:skellam-explicit-rate}
Assume the setting of Theorem~\ref{thm:skellam-shift} with the canonical calibration $e^{\varepsilon_0(n)} = c^2 n$ and $k = \lfloor\pi n\rfloor$ for some fixed $\pi \in (0,1)$. Then
\[
\Delta\bigl((P_n, Q_n), (P_\infty, Q_\infty)\bigr) \leq \frac{2c^2 + 3}{c^4 n}.
\]
In particular, the Skellam limit is attained at the explicit rate $O(n^{-1})$.
\end{corollary}

\begin{proof}
Let $r_n := \pi n - \lfloor\pi n\rfloor \in [0,1)$. Under the canonical calibration, $\delta_n = (1 + c^2 n)^{-1}$, so
\[
n\delta_n^2 = \frac{n}{(1 + c^2 n)^2} \leq \frac{1}{c^4 n}.
\]
Also
\[
\lambda_{0,n} - \lambda_0 = \frac{n - k}{1 + c^2 n} - \frac{1 - \pi}{c^2} = \frac{c^2 r_n - (1-\pi)}{c^2(1 + c^2 n)},
\]
hence $|\lambda_{0,n} - \lambda_0| \leq (c^2 + 1)/(c^4 n)$. Similarly, $|\lambda_{1,n} - \lambda_1| \leq (c^2 + 1)/(c^4 n)$. Therefore \eqref{eq:TVPn_Sk} gives
\[
\mathrm{TV}(P_n, P_\infty) \leq \frac{1}{c^4 n} + \frac{2(c^2 + 1)}{c^4 n} = \frac{2c^2 + 3}{c^4 n}.
\]
For the alternative parameters, $|\lambda'_{0,n} - \lambda_0|$ and $|\lambda'_{1,n} - \lambda_1|$ are again bounded by $(c^2+1)/(c^4 n)$, so \eqref{eq:TVQn_Sk} yields
the same bound for $\mathrm{TV}(Q_n, Q_\infty)$. Apply \eqref{eq:LeCam_bound}.
\end{proof}

\begin{proposition}[Sharp $n^{-1}$ total-variation rate in the canonical Skellam regime]
\label{prop:skellam-tv-rate}
Assume the setting of Theorem~\ref{thm:skellam-shift}, the canonical calibration $e^{\varepsilon_0(n)}=c^2 n$, and
$k_n=\lfloor \pi n\rfloor$ (more generally, the same proof works whenever $k_n-\pi n=O(1)$). Then,
for all sufficiently large $n$,
\[
\frac{e^{-1/c^2}}{4c^4}\,\frac1n
\le \mathrm{TV}(P_n,P_\infty)
\le \frac{2c^2+3}{c^4}\,\frac1n.
\]
In particular,
\[
\mathrm{TV}(P_n,P_\infty)=\Theta(n^{-1}).
\]
The same conclusion holds for $\mathrm{TV}(Q_n,Q_\infty)$.
\end{proposition}

\begin{proof}
The upper bound is Corollary~\ref{cor:skellam-explicit-rate}. For the lower bound, write
\[
\alpha_n:=k_n-\pi n,
\]
so that $\alpha_n=O(1)$. Let
\[
G_n(z):=\mathbb E[z^{D_{n,k_n}}],
\qquad
G_\infty(z):=\mathbb E[z^{D}]
=\exp\bigl(\lambda_0(z-1)+\lambda_1(z^{-1}-1)\bigr),
\]
where $D\sim\mathrm{Skellam}(\lambda_0,\lambda_1)$ with
$\lambda_0=(1-\pi)c^{-2}$ and $\lambda_1=\pi c^{-2}$. Under the canonical calibration, with
$u_n:=(c^2 n)^{-1}$ and $\delta_n=u_n/(1+u_n)$,
\[
G_n(z)=\left(\frac{1+u_n z}{1+u_n}\right)^{n-k_n}
\left(\frac{1+u_n z^{-1}}{1+u_n}\right)^{k_n}.
\]
Now evaluate at $z=i$. Using
\[
\log(1\pm iu_n)=\pm i u_n + \frac{u_n^2}{2}+O(u_n^3),
\qquad
\log(1+u_n)=u_n-\frac{u_n^2}{2}+O(u_n^3),
\]
we obtain
\[
\log G_n(i)
=\log G_\infty(i)+\frac1n\left(\frac1{c^4}-\frac{2i\alpha_n}{c^2}\right)+O(n^{-2}),
\]
and therefore
\[
G_n(i)-G_\infty(i)
=\frac{G_\infty(i)}{n}\left(\frac1{c^4}-\frac{2i\alpha_n}{c^2}\right)+O(n^{-2}).
\]
Since $|G_\infty(i)|=e^{-1/c^2}$ and
\[
\left|\frac1{c^4}-\frac{2i\alpha_n}{c^2}\right|\ge \frac1{c^4},
\]
we get, for all sufficiently large $n$,
\[
|G_n(i)-G_\infty(i)|\ge \frac{e^{-1/c^2}}{2c^4}\,\frac1n.
\]
Finally, for any two laws $\mu,\nu$ on $\mathbb Z$ and any $|z|=1$,
\[
\bigl|\mathbb E_\mu z^X-\mathbb E_\nu z^Y\bigr|
\le \sum_{d\in\mathbb Z}|\mu\{d\}-\nu\{d\}|
=2\,\mathrm{TV}(\mu,\nu).
\]
Applying this with $z=i$ yields
\[
\mathrm{TV}(P_n,P_\infty)
\ge \frac12 |G_n(i)-G_\infty(i)|
\ge \frac{e^{-1/c^2}}{4c^4}\,\frac1n.
\]
This proves $\mathrm{TV}(P_n,P_\infty)=\Theta(n^{-1})$. The same argument applies to $Q_n$, noting that
$Q_n$ is the law of $1+A'_{n,k_n}-B'_{n,k_n}$, with the replacement $\alpha_n\mapsto \alpha_n+1$, and
that shifts preserve total variation.
\end{proof}

\begin{corollary}[Privacy curve convergence for the Skellam-shift limit]
\label{cor:skellam-privacy-curve-convergence}
Under the assumptions of Theorem~\ref{thm:skellam-shift}, for every fixed $\varepsilon \geq 0$,
\[
\delta_{Q_n\|P_n}(\varepsilon) \to \delta_{Q_\infty\|P_\infty}(\varepsilon).
\]
Under the canonical calibration $e^{\varepsilon_0(n)} = c^2 n$ and $k = \lfloor \pi n \rfloor$, the convergence is quantitative:
\[
\bigl|\delta_{Q_n\|P_n}(\varepsilon) - \delta_{Q_\infty\|P_\infty}(\varepsilon)\bigr|
\leq (1 + e^\varepsilon)\,\frac{2c^2 + 3}{c^4 n}.
\]
\end{corollary}

\begin{proof}
Apply Lemma~\ref{lem:privacy-curve-stability} with the TV bound of Corollary~\ref{cor:skellam-explicit-rate}.
\end{proof}

\subsection{Skellam pmf and limiting privacy curve}

\begin{remark}[Skellam pmf]
If $D \sim \mathrm{Skellam}(\lambda_0, \lambda_1)$ with $\lambda_0, \lambda_1 > 0$, then for every $d \in \mathbb{Z}$,
\[
P(D = d) = e^{-(\lambda_0 + \lambda_1)} \left(\frac{\lambda_0}{\lambda_1}\right)^{d/2} I_{|d|}\!\left(2\sqrt{\lambda_0 \lambda_1}\right),
\]
where $I_\nu(\cdot)$ is the modified Bessel function of the first kind. This explicit pmf is not needed for the convergence
proofs, but it is convenient for numerical evaluation of the limiting privacy curve.
\end{remark}

\begin{corollary}[Limiting privacy curve; no interior $\delta$-floor]
\label{cor:skellam-curve}
Let $D \sim \mathrm{Skellam}(\lambda_0, \lambda_1)$ and write $p_\infty(d) := P(D = d)$. Let
$Q_\infty$ be the law of $1 + D$, so $Q_\infty(d) = p_\infty(d-1)$. Then for every
$\varepsilon \geq 0$,
\begin{equation}
\delta_{Q_\infty\|P_\infty}(\varepsilon) = \sum_{d \in \mathbb{Z}} \bigl(p_\infty(d-1) - e^\varepsilon p_\infty(d)\bigr)_+. \label{eq:Sk_curve}
\end{equation}
If $\pi \in (0,1)$ then $\lambda_0, \lambda_1 > 0$ and $p_\infty(d) > 0$ for all $d \in \mathbb{Z}$, hence there is no two-sided support-mismatch floor
(contrast Proposition~\ref{prop:poisson-curve}). At boundary compositions $\pi \in \{0,1\}$ one of $\lambda_0, \lambda_1$ vanishes and the limit reduces
to the Poisson-shift experiment of Section~3.
\end{corollary}

\begin{proof}
The series form \eqref{eq:Sk_curve} is the countable-space identity \eqref{eq:NP2} in Lemma~2.1 applied on $\mathbb{Z}$ with $Q_\infty(d) = p_\infty(d-1)$. If $\lambda_0, \lambda_1 > 0$, then for each $d \in \mathbb{Z}$,
\[
p_\infty(d) = P(X - Y = d) = \sum_{m \geq \max\{0,-d\}} P(X = m+d)\, P(Y = m),
\]
where $X \sim \mathrm{Poi}(\lambda_0)$ and $Y \sim \mathrm{Poi}(\lambda_1)$ are independent. Each summand is strictly positive and at least one
term appears in the sum, hence $p_\infty(d) > 0$ for all $d$, so $P_\infty$ and $Q_\infty$ share full support on $\mathbb{Z}$ and there is
no support-mismatch floor. If $\pi \in \{0,1\}$, then one of $\lambda_0, \lambda_1$ equals 0 and $D$ becomes $\pm\mathrm{Poi}(\lambda)$ for $\lambda = c^{-2}$,
which is equivalent to the Poisson-shift experiment after an affine reparametrization.
\end{proof}

\begin{remark}[Continuity of the Skellam limit at boundary compositions]
\label{rem:skellam-boundary-continuity}
The reduction to the Poisson-shift experiment at $\pi\in\{0,1\}$ stated in Corollary~\ref{cor:skellam-curve}
is continuous in the following sense. Fix the canonical calibration and let $\pi\in(0,1)$.
The total-variation bounds \eqref{eq:TVPn_Sk}--\eqref{eq:TVQn_Sk} depend on $\pi$ only through
$\lambda_0=(1-\pi)/c^2$ and $\lambda_1=\pi/c^2$. As $\pi\to 0$, we have $\lambda_1\to 0$, hence
$Y\sim\mathrm{Poi}(\lambda_1)$ vanishes in probability and
$D=X-Y\sim\mathrm{Skellam}(\lambda_0,\lambda_1)$ converges in total variation to
$\mathrm{Poi}(\lambda_0)$. Equivalently, in the centered decomposition
$K_{n,k}=k+D_{n,k}$ and $K_{n,k+1}=k+D'_{n,k}$ from \eqref{eq:Dnk}--\eqref{eq:Dnk_alt}, one gets
$D_{n,k}\Rightarrow \mathrm{Poi}(\lambda_0)$ and $D'_{n,k}\Rightarrow 1+\mathrm{Poi}(\lambda_0)$, exactly the
Poisson-shift experiment of Section~3. As $\pi\to 1$, the symmetric statement gives
$\mathrm{Skellam}(\lambda_0,\lambda_1)\to -\mathrm{Poi}(\lambda_1)$, which is the same scalar
Poisson-shift family after the reflection $d\mapsto -d$. Thus the three-regime diagram of
Section~6 is continuous at the boundary compositions: the Skellam family interpolates between
the two Poisson-shift families without a jump in the Le Cam distance.
\end{remark}

\begin{remark}[Monotonicity of the Skellam-shift limiting curve]
Fix $\pi \in (0,1)$ and write $(P_c, Q_c)$ for the Skellam-shift experiment of Theorem~\ref{thm:skellam-shift} with $\lambda_0(c) = (1-\pi)/c^2$ and $\lambda_1(c) = \pi/c^2$.
\begin{enumerate}
\item[(i)] For fixed $c > 0$, the one-sided curve $\varepsilon \mapsto \delta_{Q_c\|P_c}(\varepsilon)$ is strictly decreasing on $[0,\infty)$. Indeed, if
$L_c(d) := Q_c(d)/P_c(d) = p_c(d-1)/p_c(d)$, then $L_c(d) \to \infty$ as $d \to \infty$, so the same argument
as in Proposition~\ref{prop:poisson-monotonicity}(i) applies.
\end{enumerate}
Monotonicity in $c$ and for the two-sided curve is established in Proposition~\ref{prop:skellam-monotonicity} below.
\end{remark}

\begin{proposition}[Monotonicity in the signal-to-noise parameter for the Skellam-shift limiting curve]
\label{prop:skellam-monotonicity}
Fix $\pi\in(0,1)$ and write $(P_c,Q_c)$ for the Skellam-shift experiment of Theorem~\ref{thm:skellam-shift} with
\[
\lambda_0(c)=\frac{1-\pi}{c^2},\qquad \lambda_1(c)=\frac{\pi}{c^2}.
\]
Then for every $\varepsilon\ge 0$ and every $0<c_1<c_2$,
\[
\delta_{Q_{c_1}\|P_{c_1}}(\varepsilon)\le \delta_{Q_{c_2}\|P_{c_2}}(\varepsilon),
\qquad
\delta_{P_{c_1}\|Q_{c_1}}(\varepsilon)\le \delta_{P_{c_2}\|Q_{c_2}}(\varepsilon).
\]
Consequently,
\[
\delta_{\mathrm{two},c_1}(\varepsilon)\le \delta_{\mathrm{two},c_2}(\varepsilon).
\]
Equivalently, smaller $c$ (more noise) improves privacy.
\end{proposition}

\begin{proof}
Set $t_i:=c_i^{-2}$, so $t_1>t_2$. For $i=1,2$, let
\[
D_{t_i}=X_i-Y_i,
\qquad
X_i\sim \mathrm{Poi}((1-\pi)t_i),
\qquad
Y_i\sim \mathrm{Poi}(\pi t_i),
\]
with $X_i$ and $Y_i$ independent. Thus $D_{t_i}\sim \mathrm{Skellam}((1-\pi)t_i,\pi t_i)$, so
$P_{c_i}=\mathcal L(D_{t_i})$ and $Q_{c_i}=\mathcal L(1+D_{t_i})$.

To couple the two parameters explicitly, let
\[
U\sim \mathrm{Poi}((1-\pi)(t_1-t_2)),
\qquad
V\sim \mathrm{Poi}(\pi(t_1-t_2)),
\]
independent of $(X_2,Y_2)$, and define
\[
R:=U-V.
\]
Then
\[
R\sim \mathrm{Skellam}((1-\pi)(t_1-t_2),\pi(t_1-t_2)),
\qquad
R\ \text{is independent of } D_{t_2}.
\]
Now set $X_1:=X_2+U$ and $Y_1:=Y_2+V$. By additivity of independent Poisson variables,
\[
X_1\sim \mathrm{Poi}((1-\pi)t_1),
\qquad
Y_1\sim \mathrm{Poi}(\pi t_1),
\]
and therefore
\[
D_{t_1}=X_1-Y_1=(X_2-Y_2)+(U-V)=D_{t_2}+R.
\]
So $(P_{c_1},Q_{c_1})$ is obtained from $(P_{c_2},Q_{c_2})$ by adding the same independent
Skellam noise to both hypotheses. Equivalently, if
\[
K(d,A):=\mathbb P(d+R\in A),\qquad d\in\mathbb Z,\ A\subseteq \mathbb Z,
\]
then $K$ is a Markov kernel on $\mathbb Z$ and
\[
P_{c_1}=P_{c_2}K,
\qquad
Q_{c_1}=Q_{c_2}K.
\]

Let
\[
p_i(d):=P_{c_i}\{d\},
\qquad
q_i(d):=Q_{c_i}\{d\}=p_i(d-1),
\qquad
\rho(r):=\mathbb P(R=r).
\]
The kernel representation gives the convolution identities
\[
p_1(d)=\sum_{r\in\mathbb Z} p_2(d-r)\rho(r),
\qquad
q_1(d)=\sum_{r\in\mathbb Z} q_2(d-r)\rho(r)
      =\sum_{r\in\mathbb Z} p_2(d-1-r)\rho(r).
\]

We now write out data processing for privacy curves using the countable-space series formula
(compare \eqref{eq:NP2}; in the present one-dimensional shift setting this is exactly the Skellam
series \eqref{eq:Sk_curve}). For the forward curve,
\begin{align*}
\delta_{Q_{c_1}\|P_{c_1}}(\varepsilon)
&= \sum_{d\in\mathbb Z}\bigl(q_1(d)-e^{\varepsilon}p_1(d)\bigr)_+ \\
&= \sum_{d\in\mathbb Z}\Bigl(\sum_{r\in\mathbb Z}\rho(r)
   \bigl[p_2(d-1-r)-e^{\varepsilon}p_2(d-r)\bigr]\Bigr)_+ \\
&\le \sum_{d\in\mathbb Z}\sum_{r\in\mathbb Z}\rho(r)
   \bigl(p_2(d-1-r)-e^{\varepsilon}p_2(d-r)\bigr)_+
   \qquad\text{since }\bigl(\sum_r a_r\bigr)_+\le \sum_r (a_r)_+ \\
&= \sum_{r\in\mathbb Z}\rho(r)
   \sum_{d\in\mathbb Z}\bigl(p_2(d-1-r)-e^{\varepsilon}p_2(d-r)\bigr)_+ \\
&= \sum_{r\in\mathbb Z}\rho(r)
   \sum_{u\in\mathbb Z}\bigl(p_2(u-1)-e^{\varepsilon}p_2(u)\bigr)_+
   \qquad (u=d-r) \\
&= \sum_{r\in\mathbb Z}\rho(r)\,\delta_{Q_{c_2}\|P_{c_2}}(\varepsilon)
 = \delta_{Q_{c_2}\|P_{c_2}}(\varepsilon),
\end{align*}
as claimed.

The reverse curve is handled in exactly the same way:
\begin{align*}
\delta_{P_{c_1}\|Q_{c_1}}(\varepsilon)
&= \sum_{d\in\mathbb Z}\bigl(p_1(d)-e^{\varepsilon}q_1(d)\bigr)_+ \\
&= \sum_{d\in\mathbb Z}\Bigl(\sum_{r\in\mathbb Z}\rho(r)
   \bigl[p_2(d-r)-e^{\varepsilon}p_2(d-1-r)\bigr]\Bigr)_+ \\
&\le \sum_{d\in\mathbb Z}\sum_{r\in\mathbb Z}\rho(r)
   \bigl(p_2(d-r)-e^{\varepsilon}p_2(d-1-r)\bigr)_+ \\
&= \sum_{r\in\mathbb Z}\rho(r)
   \sum_{u\in\mathbb Z}\bigl(p_2(u)-e^{\varepsilon}p_2(u-1)\bigr)_+
   \qquad (u=d-r) \\
&= \sum_{r\in\mathbb Z}\rho(r)\,\delta_{P_{c_2}\|Q_{c_2}}(\varepsilon)
 = \delta_{P_{c_2}\|Q_{c_2}}(\varepsilon).
\end{align*}
Therefore
\[
\delta_{\mathrm{two},c_1}(\varepsilon)
= \max\{\delta_{Q_{c_1}\|P_{c_1}}(\varepsilon),\delta_{P_{c_1}\|Q_{c_1}}(\varepsilon)\}
\le
\max\{\delta_{Q_{c_2}\|P_{c_2}}(\varepsilon),\delta_{P_{c_2}\|Q_{c_2}}(\varepsilon)\}
= \delta_{\mathrm{two},c_2}(\varepsilon).
\]
This is exactly the data-processing inequality for privacy curves, written out explicitly for the
additive-noise kernel $K$.
\end{proof}


\section{General alphabets: multivariate compound-Poisson / independent Poisson vector limit}

We now allow a general finite output alphabet $\mathcal{Y}$ and consider $n$-dependent local randomizers $W^{(n)}$
that become increasingly concentrated on dominant outputs as $n \to \infty$.

\subsection{Sparse-error critical regime}

\begin{definition}[Sparse-error critical regime]
\label{def:sparse-error}
We say that $W^{(n)}$ is in the \emph{sparse-error critical regime} with dominant outputs $(y_0, y_1)$ if $y_0 \neq y_1$ and there exist nonnegative intensities $\{\alpha_0(y)\}_{y \neq y_0}$
and $\{\alpha_1(y)\}_{y \neq y_1}$ such that, as $n \to \infty$,
\[
n W^{(n)}_0(y) \to \alpha_0(y) \quad (y \neq y_0), \qquad n W^{(n)}_1(y) \to \alpha_1(y) \quad (y \neq y_1),
\]
and simultaneously $W^{(n)}_0(y_0) \to 1$ and $W^{(n)}_1(y_1) \to 1$.
\end{definition}

\begin{remark}[Limits of the ``single dominant output'' assumption]
Definition~\ref{def:sparse-error} assumes a single dominant output under each input. This covers binary RR (with $\mathcal{Y} = \{0,1\}$) and more generally
nearly-deterministic channels where all non-dominant outputs occur with probability $\Theta(1/n)$. It
does not directly cover channels with multiple dominant outputs whose probabilities stay bounded
away from 0 as $n \to \infty$ (e.g.\ multiway randomized response with several outputs of comparable
mass). Extending the critical non-Gaussian limit to such multi-dominant families naturally leads to
a L\'{e}vy--Khintchine type decomposition in which Gaussian and compound-Poisson components may
coexist.
\end{remark}

\subsection{Two dominant outputs: the first Gaussian / compound-Poisson decomposition}

The single-dominant regime of Definition~\ref{def:sparse-error} is the special case where all $O(1)$ mass under each
input collapses onto one atom. If, instead, the $O(1)$ mass is split between two dominant outputs,
then the dominant block fluctuates on the $\sqrt n$ scale and is therefore Gaussian, while the truly
non-dominant outputs still occur with probability $\Theta(1/n)$ and generate an $O(1)$ jump field.
The resulting boundary object is the first hybrid Gaussian / compound-Poisson limit, i.e. the first
nontrivial instance of a L\'evy--Khintchine type decomposition in the present shuffle setting.

The next proposition identifies the weak hybrid Gaussian / compound-Poisson limit in the 
two-dominant regime. While the weak convergence established here does not by itself yield 
total-variation or Le Cam convergence (see Remark~\ref{rem:convergence-modes}), 
Appendix~\ref{app:hybrid-privacy} shows that the privacy curves of the full experiment 
nevertheless converge, with an explicit $O(n^{-1/2})$ rate for interior compositions 
$\pi\in(0,1)$.

\begin{definition}[Two-dominant sparse-error critical regime]
\label{def:two-dominant-sparse-error}
Fix, for each $b\in\{0,1\}$, a dominant pair
\[
D_b:=\{y_{ba},y_{bb}\}\subseteq \mathcal Y,
\qquad
p_b\in(0,1),
\]
with $y_{ba}\neq y_{bb}$. We say that the triangular array of local randomizers $W^{(n)}$ is in the
\emph{two-dominant sparse-error critical regime} if, for each $b\in\{0,1\}$,
\[
W_b^{(n)}(y_{ba})=p_b+O(n^{-1}),
\qquad
W_b^{(n)}(y_{bb})=1-p_b+O(n^{-1}),
\]
and, for every $y\notin D_b$,
\[
nW_b^{(n)}(y)\to \alpha_b(y)\in[0,\infty).
\]
Equivalently, the total non-dominant mass is of order $n^{-1}$, whereas the split between the two dominant
outputs is of order one.
\end{definition}

For the statements below we impose the simplifying disjointness assumption $D_0\cap D_1=\varnothing$. The overlapping case requires additional bookkeeping and is not pursued here.

For $b\in\{0,1\}$, write
\[
\mu_b:=p_b e_{y_{ba}}+(1-p_b)e_{y_{bb}}\in \mathbb R^{\mathcal Y},
\qquad
g_b:=e_{y_{ba}}-e_{y_{bb}}\in \mathbb R^{\mathcal Y},
\]
and let
\[
M:=\mathrm{span}\{g_0,g_1\}\subseteq \mathbb R^{\mathcal Y}.
\]
Let $\Pi_G$ be the orthogonal projection onto $M$, and let $\Pi_J:=I-\Pi_G$ be the orthogonal projection
onto $M^\perp$. Thus $\Pi_G$ extracts the dominant $\sqrt n$-fluctuations, whereas $\Pi_J$ kills the
pure dominant block and retains only the $O(1)$ jump component.

Given a composition sequence $k_n$ with $\pi_n:=k_n/n\to\pi\in[0,1]$, define the centered released histogram
\[
\widehat H_{n,k_n}:=N_{n,k_n}-(n-k_n)\mu_0-k_n\mu_1\in \mathbb R^{\mathcal Y},
\]
and the hybrid normalized statistic
\[
S_n:=\Bigl(n^{-1/2}\Pi_G\widehat H_{n,k_n},\, \Pi_J\widehat H_{n,k_n}\Bigr)
\in M\times M^\perp.
\]
Under the neighboring alternative $T_{n,k_n+1}$ we use the \emph{same} centering by $(n-k_n)\mu_0+k_n\mu_1$.

\begin{proposition}[Hybrid Gaussian / compound-Poisson limit]
\label{prop:hybrid-gaussian-compound-poisson}
Assume the two-dominant sparse-error critical regime above. Also assume
$D_0\cap D_1=\varnothing$. Finally suppose $\pi_n\to\pi\in[0,1]$. Define the Gaussian
covariance operator
\[
\Sigma:=(1-\pi)p_0(1-p_0)\,g_0g_0^{\top}+\pi p_1(1-p_1)\,g_1g_1^{\top}
\qquad \text{on } M,
\]
and, for each $y\notin D_b$, the jump vector
\[
j_{b,y}:=\Pi_J(e_y-\mu_b)\in M^\perp.
\]
Let $\nu$ be the finite measure on $M^\perp$ given by
\[
\nu
:=\sum_{y\notin D_0}(1-\pi)\alpha_0(y)\,\delta_{j_{0,y}}
 +\sum_{y\notin D_1}\pi\alpha_1(y)\,\delta_{j_{1,y}}.
\]
Let $G\sim N(0,\Sigma)$ and let $J$ be an independent compound-Poisson random vector with L\'evy measure
$\nu$, equivalently
\[
J\overset{d}=\sum_{y\notin D_0}U_y j_{0,y}+\sum_{y\notin D_1}V_y j_{1,y},
\]
where the coordinates are independent and
\[
U_y\sim \mathrm{Poi}((1-\pi)\alpha_0(y)),
\qquad
V_y\sim \mathrm{Poi}(\pi\alpha_1(y)).
\]
Set
\[
\Delta:=\Pi_J(\mu_1-\mu_0)\in M^\perp.
\]
Let $P_n$ be the law of $S_n$ under $T_{n,k_n}$, and let $Q_n$ be the law of the same statistic $S_n$
under $T_{n,k_n+1}$. Then
\[
P_n\Longrightarrow P_\infty:=\mathcal L(G,J),
\qquad
Q_n\Longrightarrow Q_\infty:=\mathcal L(G,J+\Delta).
\]
Equivalently, for $u\in M$ and $v\in M^\perp$,
\[
\mathbb E\exp\bigl(i\langle u,G\rangle+i\langle v,J\rangle\bigr)
=\exp\!\left(-\frac12\langle u,\Sigma u\rangle
+\int_{M^\perp}(e^{i\langle v,z\rangle}-1)\,\nu(dz)\right),
\]
whereas under the neighboring alternative the limiting characteristic function is multiplied by the shift factor
$e^{i\langle v,\Delta\rangle}$.

In particular, the limiting normalized experiment factors into a Gaussian component on the dominant block and an
independent compound-Poisson component on the rare-jump block.
\end{proposition}

\begin{proof}
Let $m_{0,n}:=n-k_n$ and $m_{1,n}:=k_n$. Under $T_{n,k_n}$ the released histogram is
\[
N_{n,k_n}=\sum_{i=1}^{m_{0,n}} e_{Y^{(0)}_{i,n}}+\sum_{j=1}^{m_{1,n}} e_{Y^{(1)}_{j,n}},
\]
where $Y^{(b)}_{\ell,n}\sim W_b^{(n)}$ independently. Hence
\[
\widehat H_{n,k_n}
=\sum_{i=1}^{m_{0,n}}X^{(0)}_{i,n}+\sum_{j=1}^{m_{1,n}}X^{(1)}_{j,n},
\qquad
X^{(b)}_{\ell,n}:=e_{Y^{(b)}_{\ell,n}}-\mu_b.
\]
So $S_n$ is the sum of independent triangular-array increments
\[
\xi^{(b)}_{\ell,n}:=\Bigl(n^{-1/2}\Pi_G X^{(b)}_{\ell,n},\,\Pi_J X^{(b)}_{\ell,n}\Bigr)
\in M\times M^\perp.
\]

Fix $u\in M$ and $v\in M^\perp$, and define the one-user characteristic factor
\[
\phi_{b,n}(u,v)
:=\mathbb E\exp\!\left(i\big\langle u,n^{-1/2}\Pi_G X^{(b)}_{1,n}\big\rangle
+i\big\langle v,\Pi_J X^{(b)}_{1,n}\big\rangle\right).
\]
We now expand $\phi_{b,n}(u,v)$.

For the two dominant outputs, since $\Pi_J$ annihilates $M$ and
\[
\Pi_G(e_{y_{ba}}-\mu_b)=(1-p_b)g_b,
\qquad
\Pi_G(e_{y_{bb}}-\mu_b)=-p_b g_b,
\]
we get, using $W_b^{(n)}(y_{ba})=p_b+O(n^{-1})$, $W_b^{(n)}(y_{bb})=1-p_b+O(n^{-1})$, and
\[
W_b^{(n)}(y_{ba})+W_b^{(n)}(y_{bb})=1-\sum_{y\notin D_b}W_b^{(n)}(y),
\]
that
\begin{align*}
&\,W_b^{(n)}(y_{ba})
\exp\!\left(i\frac{1-p_b}{\sqrt n}\langle u,g_b\rangle\right)
+W_b^{(n)}(y_{bb})
\exp\!\left(-i\frac{p_b}{\sqrt n}\langle u,g_b\rangle\right) \\
&=1-\sum_{y\notin D_b}W_b^{(n)}(y)
  -\frac{p_b(1-p_b)}{2n}\langle u,g_b\rangle^2+o(n^{-1}) \\
&=1-\frac1n\sum_{y\notin D_b}\alpha_b(y)
  -\frac{p_b(1-p_b)}{2n}\langle u,g_b\rangle^2+o(n^{-1}).
\end{align*}
The linear term cancels because
\[
p_b(1-p_b)-(1-p_b)p_b=0,
\]
and the $O(n^{-1})$ perturbations in the dominant probabilities affect only the constant term above.

For a non-dominant output $y\notin D_b$, the event has probability
$W_b^{(n)}(y)=\alpha_b(y)n^{-1}+o(n^{-1})$, and its contribution is
\[
\exp\!\left(i\big\langle u,n^{-1/2}\Pi_G(e_y-\mu_b)\big\rangle
+i\big\langle v,j_{b,y}\big\rangle\right)
=e^{i\langle v,j_{b,y}\rangle}\bigl(1+O(n^{-1/2})\bigr).
\]
After multiplication by $W_b^{(n)}(y)=O(n^{-1})$, the $u$-dependence of these rare terms is only
$O(n^{-3/2})$, hence negligible on the logarithmic scale. Summing over the finitely many non-dominant
outputs yields
\[
\sum_{y\notin D_b}W_b^{(n)}(y)
\exp\!\left(i\big\langle u,n^{-1/2}\Pi_G(e_y-\mu_b)\big\rangle
+i\big\langle v,j_{b,y}\big\rangle\right)
=\frac1n\sum_{y\notin D_b}\alpha_b(y)e^{i\langle v,j_{b,y}\rangle}+o(n^{-1}).
\]
Combining dominant and non-dominant contributions, and using
$\sum_{y\notin D_b}\alpha_b(y)<\infty$ because $\mathcal Y$ is finite, we obtain
\[
\phi_{b,n}(u,v)
=1-\frac{p_b(1-p_b)}{2n}\langle u,g_b\rangle^2
+\frac1n\sum_{y\notin D_b}\alpha_b(y)\bigl(e^{i\langle v,j_{b,y}\rangle}-1\bigr)
+o(n^{-1}).
\]
Therefore
\[
\log\phi_{b,n}(u,v)
=-\frac{p_b(1-p_b)}{2n}\langle u,g_b\rangle^2
+\frac1n\sum_{y\notin D_b}\alpha_b(y)\bigl(e^{i\langle v,j_{b,y}\rangle}-1\bigr)
+o(n^{-1}),
\]
since $\phi_{b,n}(u,v)=1+O(n^{-1})$.

Because the two input groups are independent,
\[
\log \mathbb E_{T_{n,k_n}}\exp\bigl(i\langle u,S_n^{(G)}\rangle+i\langle v,S_n^{(J)}\rangle\bigr)
=m_{0,n}\log\phi_{0,n}(u,v)+m_{1,n}\log\phi_{1,n}(u,v),
\]
where $S_n=(S_n^{(G)},S_n^{(J)})$. Using $m_{0,n}/n\to 1-\pi$ and $m_{1,n}/n\to\pi$, we conclude that
\begin{align*}
&\log \mathbb E_{T_{n,k_n}}\exp\bigl(i\langle u,S_n^{(G)}\rangle+i\langle v,S_n^{(J)}\rangle\bigr) \\
&\to
-\frac12\Bigl((1-\pi)p_0(1-p_0)\langle u,g_0\rangle^2
+\pi p_1(1-p_1)\langle u,g_1\rangle^2\Bigr)
+\int_{M^\perp}(e^{i\langle v,z\rangle}-1)\,\nu(dz) \\
&=
-\frac12\langle u,\Sigma u\rangle
+\int_{M^\perp}(e^{i\langle v,z\rangle}-1)\,\nu(dz).
\end{align*}
This is exactly the characteristic exponent of an independent pair $(G,J)$ with
$G\sim N(0,\Sigma)$ and $J$ compound-Poisson with L\'evy measure $\nu$. By L\'evy's continuity theorem,
$P_n\Longrightarrow\mathcal L(G,J)$.

For the neighboring alternative $T_{n,k_n+1}$, write
\[
\widehat H_{n,k_n}^{\mathrm{alt}}
=\Bigl(N_{n,k_n+1}-(n-k_n-1)\mu_0-(k_n+1)\mu_1\Bigr)+(\mu_1-\mu_0).
\]
Hence
\[
S_n^{\mathrm{alt}}
=\widetilde S_n+
\Bigl(n^{-1/2}\Pi_G(\mu_1-\mu_0),\,\Pi_J(\mu_1-\mu_0)\Bigr),
\]
where $\widetilde S_n$ is the hybrid statistic built from the true group counts $n-k_n-1$ and $k_n+1$.
The same characteristic-function computation as above shows
\[
\widetilde S_n\Longrightarrow (G,J),
\]
because replacing $m_{0,n},m_{1,n}$ by $m_{0,n}-1,m_{1,n}+1$ changes the exponent only by $o(1)$. Also,
\[
n^{-1/2}\Pi_G(\mu_1-\mu_0)\to 0,
\qquad
\Pi_J(\mu_1-\mu_0)=\Delta.
\]
Therefore
\[
Q_n\Longrightarrow \mathcal L(G,J+\Delta),
\]
as claimed.
\end{proof}

\begin{remark}[L\'evy--Khintchine interpretation and the next step]
The limiting exponent obtained in the proof is
\[
-\frac12\langle u,\Sigma u\rangle
+\int_{M^\perp}(e^{i\langle v,z\rangle}-1)\,\nu(dz),
\]
namely a quadratic Gaussian part plus a finite jump measure. This is exactly the first L\'evy--Khintchine
layer for shuffle limits: the dominant $O(1)$ mass contributes a Brownian/Gaussian block, while the
$1/n$-rare deviations contribute a compound-Poisson jump block.

In the present neighboring calibration the Gaussian factor is common to both hypotheses, because replacing
one $0$-user by one $1$-user changes the dominant block only by an $O(1)$ amount, hence by $o(1)$ after
$\sqrt n$ normalization. If one simultaneously introduces an $n^{-1/2}$ local perturbation in the dominant
block, then the same computation yields a genuine hybrid shift experiment
\[
(G,J)
\quad\text{versus}\quad
(G+h,\,J+\Delta),
\]
with a nonzero Gaussian shift $h\in M$ and the compound-Poisson shift $\Delta\in M^\perp$. That is the
natural next step toward a full L\'evy--Khintchine universality theorem for shuffle privacy.
\end{remark}

\begin{remark}[Modes of convergence: TV versus weak; consequences for Le Cam distance]
\label{rem:convergence-modes}
It is essential to distinguish the convergence mode in Proposition~\ref{prop:hybrid-gaussian-compound-poisson}
from that in Theorems~\ref{thm:poisson-shift}, \ref{thm:skellam-shift}, and \ref{thm:multi-poisson-shift}.

\begin{enumerate}
\item[(i)] \textbf{Theorems~\ref{thm:poisson-shift}, \ref{thm:skellam-shift}, \ref{thm:multi-poisson-shift}}
establish convergence of $P_n \to P_\infty$ and $Q_n \to Q_\infty$ in \emph{total variation},
with explicit $O(n^{-1})$ rates. Since total variation controls Le Cam distance via Lemma~2.5,
these theorems yield genuine experiment-level convergence
$\Delta\bigl((P_n,Q_n),(P_\infty,Q_\infty)\bigr)\to 0$ with quantitative bounds.

\item[(ii)] \textbf{Proposition~\ref{prop:hybrid-gaussian-compound-poisson}} establishes only
\emph{weak convergence} $P_n\Rightarrow P_\infty$ and $Q_n\Rightarrow Q_\infty$ of the hybrid
statistic $S_n=(n^{-1/2}\Pi_G\widehat H,\,\Pi_J\widehat H)$. This is qualitatively different:
weak convergence does \emph{not} in general imply total variation convergence (the limiting laws
$P_\infty,Q_\infty$ have a Gaussian component on the continuous space $M$, whereas for each $n$
the Gaussian coordinate $S_n^{(G)}=n^{-1/2}\Pi_G\widehat H_{n,k_n}$ is supported on a countable set.
Any non-degenerate Gaussian law on $M$ is non-atomic, so if $G\sim N(0,\Sigma)$ is non-degenerate then
$\mathrm{TV}(\mathcal L(S_n^{(G)}),\mathcal L(G))=1$ for every $n$.) Accordingly,
Proposition~\ref{prop:hybrid-gaussian-compound-poisson} does \emph{not} yield a bound on
$\Delta\bigl((P_n,Q_n),(P_\infty,Q_\infty)\bigr)$ via Lemma~2.5, and should be understood as a
\emph{structural weak-limit result}: it characterizes the limiting \emph{shape} of the experiment
(Gaussian bulk plus compound-Poisson jump field) rather than establishing a Le Cam equivalence
statement.

To obtain a genuine Le Cam convergence result for the two-dominant regime one may either
(a)~project onto the jump component $\Pi_J S_n$, or (b)~show that the Gaussian factor is 
asymptotically common to both hypotheses and hence does not affect privacy curves.
Route~(a) is immediate from Theorem~\ref{thm:multi-poisson-shift}
applied to the projected experiment and is recorded formally in
Corollary~\ref{cor:projected-jump-lecam} in Section~5.3.
Route~(b) is carried out in Appendix~\ref{app:hybrid-privacy}: 
Lemma~\ref{lem:interior-conditional-smoothing} establishes an $O(n^{-1/2})$ 
conditional-smoothing bound for interior compositions $\pi\in(0,1)$, and 
Corollary~\ref{cor:full-hybrid-privacy-convergence} deduces privacy-curve convergence 
for the full hybrid experiment. The boundary compositions $\pi\in\{0,1\}$ are covered 
by the Poisson-shift results of Sections~3--4.
\end{enumerate}
\end{remark}


\begin{definition}[Limit compound-Poisson law]
\label{def:limit-compound-poisson}
Let $\{U_y\}_{y \neq y_0}$ be independent with $U_y \sim \mathrm{Poi}((1-\pi)\alpha_0(y))$, and let $\{V_y\}_{y \neq y_1}$
be independent with $V_y \sim \mathrm{Poi}(\pi\alpha_1(y))$, independent of the $U$'s. Define
the compound-Poisson random vector
\begin{equation}
H_\infty := \sum_{y \neq y_0} U_y(e_y - e_{y_0}) + \sum_{y \neq y_1} V_y(e_y - e_{y_1}) \in \mathbb{Z}^{\mathcal{Y}}. \label{eq:Hinf}
\end{equation}
\end{definition}

\subsection{Critical multivariate Poisson-shift limit}

Fix a composition sequence $k = k(n)$ with $\pi_n := k/n \to \pi \in [0,1]$. Let $N_{n,k} \in \mathbb{Z}^{\mathcal{Y}}$ denote the
released histogram under $T_{n,k}$. Define the centered histogram
\begin{equation}
H_{n,k} := N_{n,k} - (n-k) e_{y_0} - k e_{y_1} \in \mathbb{Z}^{\mathcal{Y}}, \label{eq:Hnk}
\end{equation}
where $e_y$ is the standard basis vector in $\mathbb{Z}^{\mathcal{Y}}$.


\begin{theorem}[Critical multivariate Poisson-shift limit]
\label{thm:multi-poisson-shift}
Assume $\pi_n = k/n \to \pi \in [0,1]$ and the sparse-error critical regime of Definition~\ref{def:sparse-error}. Let $P_n$ be the law of $H_{n,k}$ under $T_{n,k}$ and
let $Q_n$ be the law of $H_{n,k}$ under $T_{n,k+1}$ (centered using $k$ as in \eqref{eq:Hnk}). Let $P_\infty := \mathcal{L}(H_\infty)$ and
$Q_\infty := \mathcal{L}(H_\infty + e_{y_1} - e_{y_0})$, where $H_\infty$ is defined in \eqref{eq:Hinf}. Then
\[
P_n \to P_\infty, \qquad Q_n \to Q_\infty \quad \text{in total variation.}
\]
In particular, the neighboring shuffle experiment converges in Le Cam distance to the multivariate
Poisson-shift experiment $(P_\infty, Q_\infty)$.

Moreover, define the total ``error probabilities''
\[
p_{0,n} := \sum_{y \neq y_0} W^{(n)}_0(y) = 1 - W^{(n)}_0(y_0), \qquad p_{1,n} := \sum_{y \neq y_1} W^{(n)}_1(y) = 1 - W^{(n)}_1(y_1).
\]
Then the TV convergence admits the explicit bounds
\begin{align}
\mathrm{TV}(P_n, P_\infty) &\leq (n-k)p_{0,n}(1 - e^{-p_{0,n}}) + k p_{1,n}(1 - e^{-p_{1,n}}) \nonumber\\
&\quad + \sum_{y \neq y_0} \bigl|(n-k)W^{(n)}_0(y) - (1-\pi)\alpha_0(y)\bigr| + \sum_{y \neq y_1} \bigl|k W^{(n)}_1(y) - \pi\alpha_1(y)\bigr|, \label{eq:TVPn_multi}\\
\mathrm{TV}(Q_n, Q_\infty) &\leq (n-k-1)p_{0,n}(1 - e^{-p_{0,n}}) + (k+1)p_{1,n}(1 - e^{-p_{1,n}}) \nonumber\\
&\quad + \sum_{y \neq y_0} \bigl|(n-k-1)W^{(n)}_0(y) - (1-\pi)\alpha_0(y)\bigr| + \sum_{y \neq y_1} \bigl|(k+1)W^{(n)}_1(y) - \pi\alpha_1(y)\bigr|. \label{eq:TVQn_multi}
\end{align}
Under the mild regularity conditions $|\pi_n - \pi| = O(1/n)$ and $n W^{(n)}_b(y) = \alpha_b(y) + O(1/n)$ for all $y$,
both bounds are $O(1/n)$.
\end{theorem}

\begin{proof}
The argument makes precise the heuristic ``each user makes an error with probability $O(1/n)$,
so the multiset of errors is approximately an independent rare-category Poisson vector.''

\textbf{Step 1: decompose the histogram into two independent multinomials.} Under $T_{n,k}$
there are $n-k$ users with input 0 and $k$ users with input 1. Let
\[
N^{(0)}_{n,k}(y) := \sum_{i:\, x_i=0} \mathbf{1}\{Y_i = y\}, \qquad N^{(1)}_{n,k}(y) := \sum_{i:\, x_i=1} \mathbf{1}\{Y_i = y\}.
\]
By independence across users,
\[
N^{(0)}_{n,k} \sim \mathrm{Mult}(n-k, W^{(n)}_0), \qquad N^{(1)}_{n,k} \sim \mathrm{Mult}(k, W^{(n)}_1),
\]
and the two multinomials are independent. Moreover $N_{n,k} = N^{(0)}_{n,k} + N^{(1)}_{n,k}$ and
\begin{equation}
H_{n,k} = \bigl(N^{(0)}_{n,k} - (n-k)e_{y_0}\bigr) + \bigl(N^{(1)}_{n,k} - k e_{y_1}\bigr). \label{eq:Hnk_decomp}
\end{equation}

\textbf{Step 2: isolate deviation counts from the dominant outputs.} Define the deviation vectors (dropping the dominant coordinates)
\[
D^{(0)}_{n,k} := \bigl(N^{(0)}_{n,k}(y)\bigr)_{y \neq y_0} \in \mathbb{N}^{\mathcal{Y}\setminus\{y_0\}}, \qquad D^{(1)}_{n,k} := \bigl(N^{(1)}_{n,k}(y)\bigr)_{y \neq y_1} \in \mathbb{N}^{\mathcal{Y}\setminus\{y_1\}}.
\]
Define the deterministic linear map
\begin{equation}
\Phi(u, v) := \sum_{y \neq y_0} u(y)(e_y - e_{y_0}) + \sum_{y \neq y_1} v(y)(e_y - e_{y_1}). \label{eq:Phi}
\end{equation}
Then \eqref{eq:Hnk_decomp} implies $H_{n,k} = \Phi(D^{(0)}_{n,k}, D^{(1)}_{n,k})$.

\textbf{Step 3: Poisson approximation for multinomial deviation vectors.} Let $\tilde{U}^{(n)} = (\tilde{U}^{(n)}_y)_{y \neq y_0}$
have independent coordinates $\tilde{U}^{(n)}_y \sim \mathrm{Poi}((n-k)W^{(n)}_0(y))$. A standard explicit Poisson approximation for multinomials on rare categories (Lemma~A.3) yields
\begin{equation}
\mathrm{TV}\bigl(\mathcal{L}(D^{(0)}_{n,k}), \mathcal{L}(\tilde{U}^{(n)})\bigr) \leq (n-k)p_{0,n}(1 - e^{-p_{0,n}}). \label{eq:TV_D0}
\end{equation}
Similarly, let $\tilde{V}^{(n)} = (\tilde{V}^{(n)}_y)_{y \neq y_1}$ have independent coordinates $\tilde{V}^{(n)}_y \sim \mathrm{Poi}(k W^{(n)}_1(y))$. Then
\begin{equation}
\mathrm{TV}\bigl(\mathcal{L}(D^{(1)}_{n,k}), \mathcal{L}(\tilde{V}^{(n)})\bigr) \leq k p_{1,n}(1 - e^{-p_{1,n}}). \label{eq:TV_D1}
\end{equation}
Since the two groups are independent, by Lemma~2.4,
\[
\mathrm{TV}\bigl(\mathcal{L}(D^{(0)}_{n,k}, D^{(1)}_{n,k}), \mathcal{L}(\tilde{U}^{(n)}, \tilde{V}^{(n)})\bigr) \leq (n-k)p_{0,n}(1-e^{-p_{0,n}}) + k p_{1,n}(1-e^{-p_{1,n}}).
\]

\textbf{Step 4: replace finite-$n$ Poisson means by limiting means.} Let $U = (U_y)_{y \neq y_0}$ and
$V = (V_y)_{y \neq y_1}$ be as in Definition~\ref{def:limit-compound-poisson}, i.e.\ independent Poisson coordinates with means $(1-\pi)\alpha_0(y)$
and $\pi\alpha_1(y)$. Using the one-dimensional Poisson perturbation bound Lemma~A.2 coordinatewise
and Lemma~2.4 (iterated over all coordinates), we obtain
\begin{equation}
\mathrm{TV}\bigl(\mathcal{L}(\tilde{U}^{(n)}), \mathcal{L}(U)\bigr) \leq \sum_{y \neq y_0} \bigl|(n-k)W^{(n)}_0(y) - (1-\pi)\alpha_0(y)\bigr|, \label{eq:perturb_U}
\end{equation}
and similarly
\begin{equation}
\mathrm{TV}\bigl(\mathcal{L}(\tilde{V}^{(n)}), \mathcal{L}(V)\bigr) \leq \sum_{y \neq y_1} \bigl|k W^{(n)}_1(y) - \pi\alpha_1(y)\bigr|. \label{eq:perturb_V}
\end{equation}
Since $U$ and $V$ are independent, the joint TV distance between $(\tilde{U}^{(n)}, \tilde{V}^{(n)})$ and $(U, V)$ is bounded
by the sum of \eqref{eq:perturb_U} and \eqref{eq:perturb_V} by Lemma~2.4.

\textbf{Step 5: push through the deterministic map $\Phi$.} By construction, $H_{n,k} = \Phi(D^{(0)}_{n,k}, D^{(1)}_{n,k})$
and $H_\infty = \Phi(U, V)$ (compare \eqref{eq:Phi} and \eqref{eq:Hinf}). Total variation cannot increase under the measurable
map $\Phi$ (Lemma~2.3), hence
\[
\mathrm{TV}(P_n, P_\infty) \leq \mathrm{TV}\bigl(\mathcal{L}(D^{(0)}_{n,k}, D^{(1)}_{n,k}), \mathcal{L}(U, V)\bigr).
\]
Combining Steps 3--4 by the triangle inequality yields \eqref{eq:TVPn_multi}.

\textbf{Step 6: the alternative $T_{n,k+1}$.} Under $T_{n,k+1}$ there are $n-k-1$ zeros and $k+1$ ones, so
the same argument as Steps 1--5 applies with $(n-k, k)$ replaced by $(n-k-1, k+1)$, yielding
the bound \eqref{eq:TVQn_multi} for the centered statistic with $(k+1)$-centering. The statistic $H_{n,k}$ in \eqref{eq:Hnk} centers
using $k$, hence differs by a deterministic shift: if
\[
\tilde{H}_n := N_{n,k+1} - (n-k-1)e_{y_0} - (k+1)e_{y_1},
\]
then $H_{n,k} = \tilde{H}_n + e_{y_1} - e_{y_0}$. Since shifts are measurable maps and preserve TV, the convergence of
$\tilde{H}_n$ to $H_\infty$ implies $H_{n,k}$ converges to $H_\infty + e_{y_1} - e_{y_0}$, and \eqref{eq:TVQn_multi} follows.

\textbf{Step 7: Le Cam distance.} Both experiments live on the countable space $\mathbb{Z}^{\mathcal{Y}}$, so Lemma~2.5
yields experiment-level convergence.
\end{proof}

\begin{corollary}[Le Cam convergence for the projected jump experiment]
\label{cor:projected-jump-lecam}
In the setting of Proposition~\ref{prop:hybrid-gaussian-compound-poisson}, define the projected
statistic
\[
S_n^J:=\Pi_J\widehat H_{n,k_n}\in M^\perp,
\]
and let
\[
P_n^J:=\mathcal L(S_n^J) \quad \text{under } T_{n,k_n},
\qquad
Q_n^J:=\mathcal L(S_n^J) \quad \text{under } T_{n,k_n+1},
\]
where under $T_{n,k_n+1}$ the same centering by $(n-k_n)\mu_0+k_n\mu_1$ is used, exactly as in
Proposition~\ref{prop:hybrid-gaussian-compound-poisson}. Equivalently, $(P_n^J,Q_n^J)$ is the
pushforward of $(P_n,Q_n)$ under the second-coordinate projection $(u,v)\mapsto v$ from
$M\times M^\perp$ onto $M^\perp$. Let
\[
P_\infty^J:=\mathcal L(J),
\qquad
Q_\infty^J:=\mathcal L(J+\Delta),
\]
with $J$ and $\Delta$ as in Proposition~\ref{prop:hybrid-gaussian-compound-poisson}. Then
\[
P_n^J \to P_\infty^J,
\qquad
Q_n^J \to Q_\infty^J
\quad \text{in total variation.}
\]
Consequently,
\[
\Delta\bigl((P_n^J,Q_n^J),(P_\infty^J,Q_\infty^J)\bigr)
\leq
\max\bigl\{\mathrm{TV}(P_n^J,P_\infty^J),\mathrm{TV}(Q_n^J,Q_\infty^J)\bigr\}
\to 0.
\]
Thus route~\textup{(a)} in Remark~\ref{rem:convergence-modes} yields a genuine Le Cam convergence
statement after projection onto $M^\perp$. The full hybrid experiment $(P_n,Q_n)$ of Proposition~\ref{prop:hybrid-gaussian-compound-poisson}, which retains the Gaussian factor on $M$, is not resolved by this projection argument; however, Appendix~\ref{app:hybrid-privacy} establishes privacy-curve convergence for the full experiment at interior compositions $\pi\in(0,1)$ via a conditional-smoothing route.
\end{corollary}

\begin{proof}
Let
\[
\mathcal Y_J:=\Pi_J\bigl(\{e_y:y\in\mathcal Y\}\bigr)\subset M^\perp,
\]
and for $b\in\{0,1\}$ let $\widetilde W_b^{(n)}$ be the pushforward of $W_b^{(n)}$ under the map
$y\mapsto \Pi_J e_y$. Write
\[
m_b:=\Pi_J\mu_b=\Pi_J e_{y_{ba}}=\Pi_J e_{y_{bb}}\in M^\perp.
\]
We claim that $m_0\neq m_1$. Indeed, if $m_0=m_1$, then
\[
e_{y_{0a}}-e_{y_{1a}}\in M=\mathrm{span}\{g_0,g_1\}
=\mathrm{span}\{e_{y_{0a}}-e_{y_{0b}},\,e_{y_{1a}}-e_{y_{1b}}\}.
\]
So there exist $c_0,c_1\in\mathbb R$ such that
\[
e_{y_{0a}}-e_{y_{1a}}
=c_0(e_{y_{0a}}-e_{y_{0b}})+c_1(e_{y_{1a}}-e_{y_{1b}}).
\]
Since $D_0\cap D_1=\varnothing$, comparison of the coordinates at $y_{0b}$ and $y_{1b}$ yields
$c_0=c_1=0$, which is impossible because $e_{y_{0a}}-e_{y_{1a}}\neq 0$. Hence $m_0\neq m_1$.
Moreover, for every $w\in\mathcal Y_J\setminus\{m_b\}$,
\[
n\widetilde W_b^{(n)}(w)\to \widetilde\alpha_b(w):=
\sum_{\substack{y\notin D_b\\ \Pi_J e_y=w}} \alpha_b(y),
\qquad
\widetilde W_b^{(n)}(m_b)\to 1.
\]
Thus $\widetilde W^{(n)}$ is in the sparse-error critical regime of Definition~\ref{def:sparse-error}
on the finite alphabet $\mathcal Y_J$ with dominant outputs $(m_0,m_1)$. Let $\widetilde H_{n,k_n}$
be the associated centered histogram as in \eqref{eq:Hnk}, and define the linear map
\[
L:\mathbb Z^{\mathcal Y_J}\to M^\perp,
\qquad
L(h):=\sum_{w\in\mathcal Y_J} h(w)\,w.
\]
By construction,
\[
L(\widetilde H_{n,k_n})=\Pi_J\widehat H_{n,k_n}=S_n^J.
\]
If $\widetilde H_\infty$ denotes the limit random vector from
Theorem~\ref{thm:multi-poisson-shift} applied to $\widetilde W^{(n)}$, then
\[
L(\widetilde H_\infty)\overset{d}=J,
\qquad
L(\widetilde H_\infty+e_{m_1}-e_{m_0})\overset{d}=J+(m_1-m_0)=J+\Delta,
\]
where the first identity is exactly the grouping of equal projected jump vectors. Therefore
Theorem~\ref{thm:multi-poisson-shift}, followed by contraction of total variation under the
measurable map $L$ (Lemma~2.3), gives
\[
\mathrm{TV}(P_n^J,P_\infty^J)\to 0,
\qquad
\mathrm{TV}(Q_n^J,Q_\infty^J)\to 0.
\]
More explicitly, these two total-variation terms are bounded by the right-hand sides of
\eqref{eq:TVPn_multi}--\eqref{eq:TVQn_multi} applied to the projected array $\widetilde W^{(n)}$.
The displayed Le Cam bound is then exactly Lemma~2.5.
\end{proof}


\begin{corollary}[Explicit $O(n^{-1})$ rate and privacy curve convergence for the multivariate limit]
\label{cor:multi-explicit-rate}
Assume the sparse-error critical regime of Definition~\ref{def:sparse-error} with $n W^{(n)}_b(y) = \alpha_b(y) + O(1/n)$
uniformly in $y$ and $b$, and $|\pi_n - \pi| = O(1/n)$. Then there exists a constant $C > 0$ depending
only on the $\alpha_b(y)$ and $\pi$ such that
\[
\Delta\bigl((P_n,Q_n),(P_\infty,Q_\infty)\bigr) \leq \frac{C}{n}.
\]
Moreover, for every fixed $\varepsilon \geq 0$,
\[
\bigl|\delta_{Q_n\|P_n}(\varepsilon) - \delta_{Q_\infty\|P_\infty}(\varepsilon)\bigr| \leq \frac{C(1+e^\varepsilon)}{n}.
\]
\end{corollary}

\begin{proof}
Under the stated regularity conditions, each term in \eqref{eq:TVPn_multi}--\eqref{eq:TVQn_multi} is $O(1/n)$:
the binomial-error terms satisfy $(n-k)p_{0,n}(1-e^{-p_{0,n}}) \leq (n-k)p_{0,n}^2 = O(1/n)$ since
$p_{0,n} = \sum_{y \neq y_0} W^{(n)}_0(y) = O(1/n)$, and similarly for the parameter-perturbation terms.
The explicit constant $C$ is the sum of all these $O(1/n)$ coefficients evaluated from \eqref{eq:TVPn_multi}--\eqref{eq:TVQn_multi}.
The privacy curve bound then follows from Lemma~\ref{lem:privacy-curve-stability}.
\end{proof}

\subsection{Limiting privacy curve as an explicit series}

\begin{corollary}[Limiting privacy curve as a series]
\label{cor:multi-curve-series}
Let $p_\infty(h) := P(H_\infty = h)$ and note that $Q_\infty(h) = p_\infty(h - e)$ with shift $e := e_{y_1} - e_{y_0}$. Then for every $\varepsilon \geq 0$,
\begin{equation}
\delta_{Q_\infty\|P_\infty}(\varepsilon) = \sum_{h \in \mathbb{Z}^{\mathcal{Y}}} \bigl(p_\infty(h - e) - e^\varepsilon p_\infty(h)\bigr)_+. \label{eq:multi_curve}
\end{equation}
Moreover:
\begin{enumerate}
\item[(a)] If $|\mathcal{Y}| = 2$, the series \eqref{eq:multi_curve} reduces to the Skellam series \eqref{eq:Sk_curve}.
\item[(b)] If $\pi = 0$ (respectively $\pi = 1$), then the limit histogram law is generated only by the 0-group (respectively
only by the 1-group), but the boundary experiment factors into a scalar Poisson-shift on the switched
coordinate and common independent noise on the remaining coordinates. Under the coordinate map $h \mapsto (h_y)_{y \neq y_0}$,
\[
P_\infty \simeq \mathrm{Poi}(\alpha_0(y_1)) \otimes \bigotimes_{y \in \mathcal{Y}\setminus\{y_0,y_1\}} \mathrm{Poi}(\alpha_0(y)),
\]
whereas
\[
Q_\infty \simeq \bigl(1 + \mathrm{Poi}(\alpha_0(y_1))\bigr) \otimes \bigotimes_{y \in \mathcal{Y}\setminus\{y_0,y_1\}} \mathrm{Poi}(\alpha_0(y)).
\]
Hence the boundary privacy curve is exactly the scalar Poisson-shift series \eqref{eq:Poi_curve} with parameter $\lambda = \alpha_0(y_1)$.
In particular, extra active coordinates contribute only common independent noise and do not affect the
privacy profile. The case $\pi = 1$ is symmetric with $y_0$ replacing $y_1$.
\end{enumerate}
\end{corollary}

\begin{proof}
Equation \eqref{eq:multi_curve} is the countable-space identity \eqref{eq:NP2} in Lemma~2.1 on the countable space $\mathbb{Z}^{\mathcal{Y}}$, using the
shift relationship $Q_\infty(h) = p_\infty(h - e)$.

\textbf{(a) Reduction to Skellam when $|\mathcal{Y}| = 2$.} Let $\mathcal{Y} = \{y_0, y_1\}$. Then there are no other coordinates, and
the only possible deviations are swaps between $y_0$ and $y_1$. In Definition~\ref{def:limit-compound-poisson}, the non-dominant sets are
$\mathcal{Y}\setminus\{y_0\} = \{y_1\}$ and $\mathcal{Y}\setminus\{y_1\} = \{y_0\}$, hence $U_{y_1} \sim \mathrm{Poi}((1-\pi)\alpha_0(y_1))$ and $V_{y_0} \sim \mathrm{Poi}(\pi\alpha_1(y_0))$ and
\[
H_\infty = U_{y_1}(e_{y_1} - e_{y_0}) + V_{y_0}(e_{y_0} - e_{y_1}) = (U_{y_1} - V_{y_0})(e_{y_1} - e_{y_0}).
\]
Thus $H_\infty$ is one-dimensional along $e_{y_1} - e_{y_0}$ with coefficient $U_{y_1} - V_{y_0} \sim \mathrm{Skellam}(\lambda_0, \lambda_1)$ after identifying
$\lambda_0 = (1-\pi)\alpha_0(y_1)$ and $\lambda_1 = \pi\alpha_1(y_0)$. The shift by $e = e_{y_1} - e_{y_0}$ corresponds to the $+1$ shift in the scalar
Skellam-shift experiment, so \eqref{eq:multi_curve} reduces to \eqref{eq:Sk_curve}.

\textbf{(b) Boundary case $\pi = 0$.} Then $V_y \equiv 0$ a.s.\ for all $y \neq y_1$, and
\[
H_\infty = \sum_{y \neq y_0} U_y(e_y - e_{y_0}), \qquad U_y \sim \mathrm{Poi}(\alpha_0(y)) \text{ independently.}
\]
Under the coordinate map $h \mapsto (h_y)_{y \neq y_0}$, the law $P_\infty$ is identified with the independent vector $(U_y)_{y \neq y_0}$,
hence
\[
P_\infty \simeq \mathrm{Poi}(\alpha_0(y_1)) \otimes \bigotimes_{y \in \mathcal{Y}\setminus\{y_0,y_1\}} \mathrm{Poi}(\alpha_0(y)).
\]
Under $Q_\infty$, shifting by $e = e_{y_1} - e_{y_0}$ increments only the $y_1$-coordinate, so
\[
Q_\infty \simeq \bigl(1 + \mathrm{Poi}(\alpha_0(y_1))\bigr) \otimes \bigotimes_{y \in \mathcal{Y}\setminus\{y_0,y_1\}} \mathrm{Poi}(\alpha_0(y)).
\]
Thus the boundary experiment factors into a scalar Poisson-shift and common independent noise, hence \eqref{eq:multi_curve}
reduces to \eqref{eq:Poi_curve}. The case $\pi = 1$ is symmetric.
\end{proof}

\begin{remark}[Boundary compositions are scalar for privacy]
\label{rem:boundary-scalar-privacy}
At boundary compositions $\pi \in \{0,1\}$ the released histogram can remain genuinely multivariate when several rare coordinates have positive limiting
intensities. Nevertheless, the binary limit experiment relevant for privacy factors into a scalar Poisson-shift
on the switched coordinate and common independent noise on the remaining coordinates. Therefore the
boundary privacy curve is scalar, not genuinely multivariate.
\end{remark}

\section{A three-regime synthesis under convergent macroscopic scalings}

Taken together, Part~I and the present paper suggest the following three canonical regimes in terms
of the scaling parameter $a_n = e^{\varepsilon_0}/n$ (or, more generally, $1/n$-level deviation
intensities in Definition~\ref{def:sparse-error}), provided the macroscopic parameters of the array
under study converge along the sequence.

\begin{itemize}
\item \textbf{Sub-critical: $a_n \to 0$ (Gaussian/GDP).} Standardized jump sizes vanish and a Lindeberg
condition holds. The privacy loss is asymptotically Gaussian and the neighboring experiment is
asymptotically equivalent to a Gaussian shift experiment as in Part~I~\cite{Shv26}, with GDP
calibration in the sense of~\cite{DRS22}. Proposition~\ref{prop:power-law-gaussian} below gives a
concrete power-law RR example with $e^{\varepsilon_0(n)} = n^\alpha$, $\alpha \in (0,1)$.
\item \textbf{Critical: $a_n \to c^2 \in (0,\infty)$ (Poisson/Skellam/compound-Poisson).}
Theorem~\ref{thm:poisson-shift} gives the Poisson-shift limit for $\pi = 0$,
Theorem~\ref{thm:skellam-shift} gives the Skellam-shift limit for $\pi \in (0,1)$, and
Theorem~\ref{thm:multi-poisson-shift} gives the multivariate compound-Poisson limit for general alphabets.
\item \textbf{Super-critical: $a_n \to \infty$ (no privacy).} Proposition~\ref{prop:supercritical-rr}
proves asymptotic distinguishability for the canonical RR pair, and
Corollary~\ref{cor:supercritical-compositions} extends the same conclusion to arbitrary RR compositions.
Corollary~\ref{cor:supercritical-finite} gives the analogous finite-alphabet sparse-error statement.
\end{itemize}

\begin{remark}[Need for convergent macroscopic parameters]
Even in the critical window, a full-sequence limit need not exist without convergence of the macroscopic
composition parameters. For example, in binary RR with $e^{\varepsilon_0(n)} = c^2 n$, let
\[
k_n=
\begin{cases}
0, & n \text{ even},\\
\lfloor n/2\rfloor, & n \text{ odd}.
\end{cases}
\]
Along the even subsequence Theorem~\ref{thm:poisson-shift} yields the canonical Poisson-shift limit,
whereas along the odd subsequence Theorem~\ref{thm:skellam-shift} yields the interior Skellam-shift limit.
These two limit experiments have different privacy curves: the Poisson-shift reverse curve has the
support-mismatch floor $e^{-1/c^2}$ by Proposition~\ref{prop:poisson-curve}, while the interior
Skellam limit has no such floor by Corollary~\ref{cor:skellam-curve}. Hence the full sequence has no
single privacy-curve limit. The phase-diagram statements in this section should therefore be read as
descriptions of subsequential limits under convergent macroscopic parameters.
\end{remark}

\begin{proposition}[Power-law sub-critical RR is still Gaussian]
\label{prop:power-law-gaussian}
Let $k_n\in\{0,\dots,n-1\}$ with $\pi_n:=k_n/n\to\pi\in[0,1]$, and suppose
\[
e^{\varepsilon_0(n)}=n^\alpha,
\qquad
\alpha\in(0,1).
\]
Let $P_n:=T_{n,k_n}$, $Q_n:=T_{n,k_n+1}$, $\Lambda_{n,k_n}:=\log(dQ_n/dP_n)$. Write
\[
\delta_n:=\frac{1}{1+n^\alpha},
\qquad
v_n:=n\delta_n(1-\delta_n),
\qquad
\Delta_n:=1-2\delta_n,
\qquad
h_n:=\frac{\Delta_n}{\sqrt{v_n}}.
\]
Then $h_n\sim n^{(\alpha-1)/2}\to 0$, and
\[
\frac{\Lambda_{n,k_n}+\frac12 h_n^2}{h_n}\Longrightarrow N(0,1)
\quad\text{under }P_n,
\qquad
\frac{\Lambda_{n,k_n}-\frac12 h_n^2}{h_n}\Longrightarrow N(0,1)
\quad\text{under }Q_n.
\]
\end{proposition}

\begin{proof}
\textbf{Step~1: one-summand decomposition.}
Under $T_{n,k_n}$ the dataset has $m_{0,n}:=n-k_n$ zero-input users and $m_{1,n}:=k_n$ one-input
users. Under $T_{n,k_n+1}$ one zero-input user is replaced by a one-input user.
Let $X_*\sim\mathrm{Bern}(\delta_n)$ denote the output of the replaced user under $P_n$,
and let $S$ be the sum of the remaining $n-1$ outputs, independent of $X_*$.
Then $K_n=S+X_*$ under $P_n$ and $K_n=S+Y_*$ under $Q_n$ where $Y_*\sim\mathrm{Bern}(1-\delta_n)$.
Writing $f(m):=P(S=m)$, the exact one-summand formulas are
\[
P_n\{K_n=m\}=(1-\delta_n)f(m)+\delta_n f(m-1),
\qquad
Q_n\{K_n=m\}=\delta_n f(m)+(1-\delta_n)f(m-1).
\]
Define the backward ratio $\rho(m):=f(m-1)/f(m)$ (for $f(m)>0$). Then
\begin{equation}
L_n(m)=\frac{Q_n\{K_n=m\}}{P_n\{K_n=m\}}
=\frac{\delta_n+(1-\delta_n)\rho(m)}{(1-\delta_n)+\delta_n\rho(m)}
=:g(\rho(m)),
\label{eq:LR-general-k}
\end{equation}
where $g(r):=[\delta_n+(1-\delta_n)r]/[(1-\delta_n)+\delta_n r]$.

\medskip
\textbf{Step~2: Taylor expansion of $\log g$.}
One checks
\[
g(1)=1,
\qquad
(\log g)'(1)=\Delta_n,
\qquad
(\log g)''(1)=-\Delta_n,
\]
so that for $|u|\le 1/2$,
\begin{equation}
\log g(1+u)=\Delta_n u-\frac{\Delta_n}{2}u^2+O(\Delta_n|u|^3).
\label{eq:logg-taylor}
\end{equation}

\medskip
\textbf{Step~3: Poisson-binomial ratio estimate.}
Let $\mu_S:=\mathbb E[S]$ and $v_S:=\mathrm{Var}(S)$. Since $S$ is a sum of $n-1$ independent Bernoulli
random variables (with two distinct parameters) and $v_S=v_n+O(1)\to\infty$,
the first-order Edgeworth expansion for lattice Poisson-binomial sums
(\cite{Pet75}, Chapter~VII) gives, uniformly for $|z_m|\le M$ with $z_m:=(m-\mu_S)/\sqrt{v_S}$,
\[
f(m)=\frac{1}{\sqrt{v_S}}\,\phi(z_m)\bigl[1+p_1(z_m)/\sqrt{v_S}+O(v_S^{-1})\bigr],
\]
where $\phi$ is the standard Gaussian density and $p_1$ is a polynomial depending on the
third cumulant of $S$. Substituting $z_{m-1}=z_m-1/\sqrt{v_S}$ and using
\[
\frac{\phi(z_m-1/\sqrt{v_S})}{\phi(z_m)}
=\exp\Bigl(\frac{z_m}{\sqrt{v_S}}-\frac{1}{2v_S}\Bigr)
=1+\frac{z_m}{\sqrt{v_S}}+O(v_S^{-1})
\]
on $|z_m|\le M$, together with $p_1(z_m-1/\sqrt{v_S})=p_1(z_m)+O(v_S^{-1/2})$, we obtain
\begin{equation}
\rho(m)=\frac{f(m-1)}{f(m)}=1+\frac{z_m}{\sqrt{v_S}}+O(v_S^{-1}),
\qquad\text{uniformly for }|z_m|\le M.
\label{eq:ratio-edgeworth}
\end{equation}

\medskip
\textbf{Step~4: CLT and assembly.}
Write $\mu_n:=\mathbb E_{P_n}[K_n]$ and $Z_n:=(K_n-\mu_n)/\sqrt{v_n}$.
Since $v_n\asymp n^{1-\alpha}\to\infty$ and each Bernoulli summand has variance at most $1/4$,
the Lindeberg condition is satisfied and $Z_n\Rightarrow N(0,1)$ under $P_n$.
Also $h_n=\Delta_n/\sqrt{v_n}\sim n^{(\alpha-1)/2}\to 0$.

Fix $M<\infty$. On the event $\{|Z_n|\le M\}$, the backward ratio at $K_n$ satisfies
\eqref{eq:ratio-edgeworth} (after noting $\mu_S=\mu_n+O(1)$ and $v_S=v_n+O(1)$, so
$z_{K_n}=Z_n+O(v_n^{-1/2})$):
\[
\rho(K_n)-1=\frac{Z_n}{\sqrt{v_n}}+O(v_n^{-1}).
\]
Since $\Delta_n\to 1$ and $\Delta_n/v_n=h_n^2/\Delta_n=O(h_n^2)$, substituting into
\eqref{eq:logg-taylor} with $u=\rho(K_n)-1$ gives
\begin{align*}
\Lambda_{n,k_n}
&=\Delta_n\cdot\frac{Z_n}{\sqrt{v_n}}
 -\frac{\Delta_n}{2}\cdot\frac{Z_n^2}{v_n}
 +O(h_n^2)\\
&=h_n Z_n-\frac{h_n^2}{2}Z_n^2+O(h_n^2).
\end{align*}
Therefore
\[
\frac{\Lambda_{n,k_n}+\frac12 h_n^2}{h_n}
=Z_n-\frac{h_n}{2}(Z_n^2-1)+O(h_n)
=Z_n+o_{P_n}(1)
\quad\text{on }\{|Z_n|\le M\}.
\]
Since $M$ is arbitrary and $Z_n\Rightarrow N(0,1)$, letting $M\to\infty$ yields
\[
\frac{\Lambda_{n,k_n}+\frac12 h_n^2}{h_n}\Longrightarrow N(0,1)
\quad\text{under }P_n.
\]

Under $Q_n$, the same CLT gives $Z_n':=(K_n-\mu_n-\Delta_n)/\sqrt{v_n}\Rightarrow N(0,1)$
and $Z_n=Z_n'+h_n$. Repeating the expansion with $Z_n=Z_n'+h_n$,
\[
\Lambda_{n,k_n}=h_n(Z_n'+h_n)-\frac{h_n^2}{2}(Z_n'+h_n)^2+O(h_n^2)
=h_n Z_n'+\frac12 h_n^2+o_{Q_n}(h_n),
\]
hence
\[
\frac{\Lambda_{n,k_n}-\frac12 h_n^2}{h_n}\Longrightarrow N(0,1)\quad\text{under }Q_n.
\]
\end{proof}



\subsection{Super-critical no-privacy lemma}

\begin{proposition}[Super-critical RR yields asymptotic distinguishability]
\label{prop:supercritical-rr}
Consider the canonical RR neighboring pair and suppose $a_n = e^{\varepsilon_0(n)}/n \to \infty$. Let $P_n, Q_n$ be the neighboring transcript laws (equivalently, the laws of $K_n$ in \eqref{eq:Kn_null}--\eqref{eq:Kn_alt}). Then
\[
\mathrm{TV}(P_n, Q_n) \to 1.
\]
In particular, for every fixed $\varepsilon \geq 0$, $\delta_{\mathrm{two}}(\varepsilon) \to 1$, and the likelihood ratio satisfies $L_n \to 0$ in
$P_n$-probability while $L_n \to \infty$ in $Q_n$-probability.
\end{proposition}

\begin{proof}
If $a_n \to \infty$, then $\delta_n = (1 + e^{\varepsilon_0(n)})^{-1} \sim e^{-\varepsilon_0(n)}$ and the null mean satisfies
\[
n\delta_n = \frac{n}{e^{\varepsilon_0(n)}} = \frac{1}{a_n} \to 0.
\]
Hence $K_n \to 0$ in $P_n$-probability.

Under $Q_n$, we have $K_n = S_{n-1} + B_n$ as in \eqref{eq:Kn_alt}. Here $S_{n-1} \to 0$ in probability since $(n-1)\delta_n \to 0$,
while $B_n \sim \mathrm{Bern}(1 - \delta_n) \to 1$ in probability because $\delta_n \to 0$. Thus $K_n \to 1$ in $Q_n$-probability.

Therefore the event $A_n := \{K_n = 0\}$ satisfies $P_n(A_n) \to 1$ while $Q_n(A_n) \to 0$, so
\[
\mathrm{TV}(P_n, Q_n) \geq P_n(A_n) - Q_n(A_n) \to 1.
\]
This implies $\delta_{\mathrm{two}}(0) = \mathrm{TV}(P_n, Q_n) \to 1$. Moreover, for any fixed $\varepsilon \geq 0$, taking $A = A_n$ in the
definition of $\delta_{P_n\|Q_n}(\varepsilon)$ gives
\[
\delta_{P_n\|Q_n}(\varepsilon) \geq P_n(A_n) - e^\varepsilon Q_n(A_n) \to 1,
\]
so $\delta_{\mathrm{two}}(\varepsilon) \to 1$ for every fixed $\varepsilon \geq 0$.

Finally, by \eqref{eq:LR}, $L_n(0) = e^{-\varepsilon_0(n)} \to 0$ and
\[
L_n(1) = e^{-\varepsilon_0(n)} + \frac{e^{\varepsilon_0(n)} - e^{-\varepsilon_0(n)}}{n} = e^{-\varepsilon_0(n)} + \frac{e^{\varepsilon_0(n)}}{n}\bigl(1 - e^{-2\varepsilon_0(n)}\bigr) \sim a_n \to \infty.
\]
Since $K_n \to 0$ under $P_n$ and $K_n \to 1$ under $Q_n$, this yields the claimed convergence in probability.
\end{proof}

\begin{corollary}[Super-critical RR for arbitrary compositions]
\label{cor:supercritical-compositions}
Assume shuffled binary RR with $a_n = e^{\varepsilon_0(n)}/n \to \infty$, and let $k = k(n) \in \{0, \ldots, n-1\}$ be any composition sequence. Let $P_{n,k} = T_{n,k}$ and $Q_{n,k} = T_{n,k+1}$ be the neighboring transcript laws. Then
\[
\mathrm{TV}(P_{n,k}, Q_{n,k}) \to 1.
\]
Equivalently, for every fixed $\varepsilon \geq 0$, the two-sided privacy curve satisfies $\delta^{(n,k)}_{\mathrm{two}}(\varepsilon) \to 1$. If in addition
$k/n \to \pi \in [0,1]$, then under $P_{n,k}$ one has $K_{n,k} \to k$ in probability, while under $Q_{n,k}$ one has
$K_{n,k+1} \to k+1$ in probability.
\end{corollary}

\begin{proof}
Write $K_{n,k} = k + A_{n,k} - B_{n,k}$ as in \eqref{eq:Dnk}. Since
\[
E[A_{n,k}] = (n-k)\delta_n \leq n\delta_n = \frac{1}{a_n} \to 0, \qquad E[B_{n,k}] = k\delta_n \leq n\delta_n = \frac{1}{a_n} \to 0,
\]
Markov's inequality gives $A_{n,k} \to 0$ and $B_{n,k} \to 0$ in probability. Hence $K_{n,k} \to k$ in $P_{n,k}$-probability.

Likewise, under $Q_{n,k}$, $K_{n,k+1} = k + 1 + A'_{n,k} - B'_{n,k}$ with
\[
E[A'_{n,k}] = (n-k-1)\delta_n \leq \frac{1}{a_n} \to 0, \qquad E[B'_{n,k}] = (k+1)\delta_n \leq \frac{1}{a_n} \to 0.
\]
Thus $K_{n,k+1} \to k+1$ in $Q_{n,k}$-probability.

Therefore the event $A_n := \{K = k\}$ satisfies
\[
P_{n,k}(A_n) \to 1,
\qquad
Q_{n,k}(A_n) \to 0.
\]
Hence $\mathrm{TV}(P_{n,k}, Q_{n,k}) \to 1$. The conclusion
$\delta^{(n,k)}_{\mathrm{two}}(\varepsilon) \to 1$ then follows from
Proposition~\ref{prop:supercritical-rr} by the same argument.
\end{proof}

\begin{corollary}[Finite-alphabet super-critical sparse-error regime]
\label{cor:supercritical-finite}
Suppose there exist dominant outputs $y_0 \neq y_1$ such that
\[
n\bigl(1 - W^{(n)}_0(y_0)\bigr) \to 0, \qquad n\bigl(1 - W^{(n)}_1(y_1)\bigr) \to 0.
\]
Let $k = k(n) \in \{0, \ldots, n-1\}$ be any composition sequence and define $H_{n,k}$ as in \eqref{eq:Hnk}. Then under $T_{n,k}$,
\[
H_{n,k} \to 0 \quad \text{in probability,}
\]
while under $T_{n,k+1}$ (centered using $k$ as in \eqref{eq:Hnk}),
\[
H_{n,k} \to e_{y_1} - e_{y_0} \quad \text{in probability.}
\]
Consequently, $\mathrm{TV}(P_n, Q_n) \to 1$ and $\delta_{\mathrm{two}}(\varepsilon) \to 1$ for every fixed $\varepsilon \geq 0$.
\end{corollary}

\begin{proof}
Let $p_{0,n} := 1 - W^{(n)}_0(y_0)$ and $p_{1,n} := 1 - W^{(n)}_1(y_1)$. Under $T_{n,k}$, each 0-user outputs a
non-dominant symbol with probability $p_{0,n}$ and each 1-user outputs a non-dominant symbol with
probability $p_{1,n}$. Hence, by the union bound,
\[
P_{T_{n,k}}(H_{n,k} \neq 0) \leq (n-k)p_{0,n} + k p_{1,n} \leq n p_{0,n} + n p_{1,n} \to 0.
\]
So $H_{n,k} \to 0$ in probability.

Under $T_{n,k+1}$, define the $(k+1)$-centered statistic
\[
\tilde{H}_{n,k+1} := N_{n,k+1} - (n-k-1)e_{y_0} - (k+1)e_{y_1}.
\]
The same union-bound argument gives $\tilde{H}_{n,k+1} \to 0$ in probability. But $H_{n,k} = \tilde{H}_{n,k+1} + e_{y_1} - e_{y_0}$,
so $H_{n,k} \to e_{y_1} - e_{y_0}$ in probability under $T_{n,k+1}$. The event $\{H_{n,k} = 0\}$ therefore separates the two
hypotheses with asymptotic probability 1, which implies $\mathrm{TV}(P_n, Q_n) \to 1$ and hence $\delta_{\mathrm{two}}(\varepsilon) \to 1$
for each fixed $\varepsilon$.
\end{proof}

\begin{remark}[Continuity at the regime boundaries]
The critical Poisson and Skellam families interpolate continuously between the Gaussian and no-privacy edges once the direction of $c$ is interpreted correctly.
\begin{enumerate}
\item[(i)] \textbf{Gaussian edge ($c \downarrow 0$).} Since $\lambda = c^{-2} \to \infty$, the Poisson and Skellam counts admit
a Gaussian approximation after centering and scaling. Concretely, if $J_c \sim \mathrm{Poi}(c^{-2})$ and
$Z_c := c(J_c - c^{-2})$, then under the Poisson-shift null $Z_c \Rightarrow \mathcal{N}(0,1)$ (by the CLT for Poisson, see e.g.\ \cite{vdV98}), while under the alternative $c\bigl((1 + J_c) - c^{-2}\bigr) = Z_c + c \Rightarrow \mathcal{N}(c, 1)$. Likewise, for the Skellam family with fixed $\pi \in (0,1)$, if $D_c \sim \mathrm{Skellam}((1-\pi)c^{-2}, \pi c^{-2})$ and $\hat{Z}_c := c(D_c - (1-2\pi)c^{-2})$, then $\hat{Z}_c \Rightarrow \mathcal{N}(0,1)$ under the null and $\hat{Z}_c + c \Rightarrow \mathcal{N}(c,1)$ under the alternative. Thus, as $c \downarrow 0$, the critical non-Gaussian families match the GDP curve with Gaussian shift parameter $\mu = c$ to first order, exactly as predicted by the sub-critical theory of Part~I~\cite{Shv26}.
\item[(ii)] \textbf{No-privacy edge ($c \uparrow \infty$).} Here $\lambda = c^{-2} \downarrow 0$. In the Poisson family, $P_c = \mathrm{Poi}(\lambda)$ concentrates on 0 while $Q_c = 1 + \mathrm{Poi}(\lambda)$ concentrates on 1; in the Skellam family, both Poisson means vanish and the pair again collapses to $(\delta_0, \delta_1)$. Hence for every fixed $\varepsilon \geq 0$, $\delta_{\mathrm{two},c}(\varepsilon) \to 1$.
\end{enumerate}
Thus the critical families connect continuously to the Gaussian/GDP edge as $c \downarrow 0$ and to the no-privacy edge as $c \uparrow \infty$.
\end{remark}

\section{Comparison with existing amplification bounds}
\label{sec:comparison-existing-bounds}

The generic amplification theorems of Balle et al.~\cite{BBG19} and Feldman et al.~\cite{FMT21}
are indispensable in the classical regime of many small contributions. The critical window studied
here,
\[
e^{\varepsilon_0(n)}\sim c^2 n,
\qquad
\lambda:=c^{-2},
\]
is structurally different: in binary shuffled randomized response the local flip probability satisfies
\[
\delta_n=\frac{1}{1+e^{\varepsilon_0(n)}}\sim \frac{\lambda}{n},
\qquad
n\delta_n\to \lambda,
\]
so the total number of exceptional local reports is of order one rather than order $n$. This is why
Theorem~\ref{thm:poisson-shift} yields a Poisson-shift limit and Proposition~\ref{prop:poisson-curve}
yields a non-vanishing support-mismatch floor. In particular,
\[
\delta_{\mathrm{two}}(\varepsilon)\ge e^{-1/c^2}
\qquad \text{for all } \varepsilon\ge 0
\]
in the canonical Poisson-shift limit. We also note that Rényi differential privacy~\cite{Mir17}
provides another lens on shuffle amplification~\cite{GKL21}, but is not the focus of the present
Le~Cam-theoretic analysis.

\paragraph{Balle et al.~\cite{BBG19}.}
The privacy-blanket framework of~\cite{BBG19} introduces a random blanket-user count
$M_n\sim \mathrm{Bin}(n,\gamma_n)$, where $\gamma_n$ depends on the local privacy level. For
binary randomized response in the critical window $e^{\varepsilon_0(n)}=c^2 n$, one computes
\[
\mathbb E M_n = n\gamma_n \to \frac{2}{c^2}.
\]
Thus even inside the blanket decomposition there are only $O(1)$ blanket users in the critical
window: the many-blanket-user regime on which the amplification bound of~\cite{BBG19} relies
does not materialize. More concretely, the simplified asymptotic corollary stated there requires
\[
\varepsilon_0 \le \frac{1}{2}\log\frac{n}{\log(1/\delta)},
\]
whereas in the critical window $\varepsilon_0(n)=\log n + O(1)$, so this hypothesis fails for every
fixed nontrivial $\delta\in(0,1)$. The blanket-based bound therefore yields no finite
privacy guarantee in the critical regime.

\paragraph{Feldman et al.~\cite{FMT21}.}
The exact reduction of~\cite{FMT21} expresses the shuffled neighboring pair in terms of a hidden
count $C_n\sim \mathrm{Bin}(n-1,e^{-\varepsilon_0(n)})$. Under the critical scaling,
\[
\mathbb E C_n=(n-1)e^{-\varepsilon_0(n)}\to \frac{1}{c^2},
\qquad
C_n\Longrightarrow \mathrm{Poi}(1/c^2).
\]
So the exact representation in~\cite{FMT21} already reveals that the relevant hidden count
converges to a Poisson law with constant mean. However, the closed-form amplification bound
of~\cite{FMT21} is built on an $n^{-1/2}$ amplification template, which predicts
$\delta\to 0$ as $\varepsilon\to\infty$, whereas the Poisson-shift limit of
Theorem~\ref{thm:poisson-shift} establishes a strictly positive floor
$\delta_{\mathrm{two}}(\varepsilon)\ge e^{-1/c^2}$ for all $\varepsilon\ge 0$.

\medskip

The comparison is structural. The critical window $e^{\varepsilon_0}\asymp n$ is the boundary
at which the hidden counts in blanket/clone decompositions stop diverging and instead converge
to Poisson laws with constant means. Once that happens, privacy loss is governed by a Poisson
number of macroscopic jumps, not by Gaussian fluctuations of $n$ small increments. In this
regime the blanket-based bound of~\cite{BBG19} yields no finite guarantee, while the $n^{-1/2}$
amplification template of~\cite{FMT21} misses the Poisson floor.


\section{Discussion}

\begin{tcolorbox}[keybox]
\textbf{Why Poisson (not Gaussian) at criticality.} In the critical scaling $a_n = e^{\varepsilon_0}/n \to c^2$, each user makes an ``error'' with probability $\delta_n \asymp 1/n$. The number of errors is $O(1)$ and converges to Poisson; each error causes a macroscopic jump in the centered histogram and hence in the log-likelihood ratio. This is the classical ``law of small numbers'' mechanism behind Poisson approximation.

\textbf{Relation to Takagi--Liew \cite{TL26}.} Independently, Takagi and Liew~\cite{TL26} develop an asymptotic blanket-divergence framework for shuffle privacy beyond pure LDP, including Gaussian local randomizers in the sub-critical regime where classical CLT behavior still governs the shuffled privacy loss. The present paper is complementary: it resolves the critical Poisson/Skellam/compound-Poisson frontier, where error probabilities are of order $1/n$, macroscopic jumps survive in the limit, and the classical Lindeberg conditions underlying Gaussian approximations fail.

\textbf{A unified L\'{e}vy--Khintchine limit theory.} A natural open direction is to develop a unified universality theory in which both Gaussian and Poisson components can coexist, analogous to classical L\'{e}vy--Khintchine limits for sums of independent random variables. Heuristically, this should occur when $a_n \to 0$ but so slowly that rare deviations produce an $O(1)$ compound-Poisson component while the bulk still contributes a Gaussian fluctuation.

\textbf{Choosing $\varepsilon_0(n)$ for target $(\varepsilon, \delta)$.} The phase diagram gives guidance for protocol design. Roughly: (i) if one targets moderate $\delta$ and fixed central $\varepsilon$, staying in the sub-critical regime $a_n \ll 1$ keeps the privacy loss in the Gaussian/GDP world covered by Part~I when its assumptions apply; (ii) pushing $\varepsilon_0$ to the critical boundary $\varepsilon_0 \approx \log n + O(1)$ enters the non-Gaussian regime and requires Poisson/Skellam/compound-Poisson calibration; (iii) if $\varepsilon_0$ grows faster than $\log n$ (super-critical), privacy collapses. At criticality one must also track the macroscopic composition parameters, since oscillating compositions can lead to different subsequential limits.
\end{tcolorbox}

\appendix

\section{Technical lemmas}

This appendix collects explicit bounds used in the proofs.

\begin{lemma}[Binomial $\to$ Poisson in total variation (explicit coupling)]
Let $S \sim \mathrm{Bin}(m, p)$ with $p \in (0,1)$ and let $N \sim \mathrm{Poi}(mp)$. Then
\begin{equation}
\mathrm{TV}(\mathcal{L}(S), \mathrm{Poi}(mp)) \leq mp(1 - e^{-p}) \leq mp^2. \label{eq:A1}
\end{equation}
\end{lemma}

\begin{proof}
We give an explicit coupling. Let $N_1, \ldots, N_m$ be i.i.d.\ $\mathrm{Poi}(p)$ and set $N := \sum_{i=1}^m N_i$.
Then $N \sim \mathrm{Poi}(mp)$.

For each $i$, define a Bernoulli random variable $X_i$ from $N_i$ by
\[
X_i := \begin{cases} 1, & N_i \geq 1, \\ \mathrm{Bern}(q), & N_i = 0, \end{cases} \qquad q := \frac{p - (1 - e^{-p})}{e^{-p}}.
\]
First note that $q \in [0,1]$. Indeed, since $1 - e^{-p} \leq p$, the numerator is nonnegative and hence $q \geq 0$.
Also
\[
q = \frac{p - 1 + e^{-p}}{e^{-p}} = 1 + (p-1)e^p = 1 - (1-p)e^p \leq 1,
\]
and the lower bound follows from $\log(1/(1-p)) \geq p$.

By construction,
\[
P(X_i = 1) = P(N_i \geq 1) + P(N_i = 0)\, q = (1 - e^{-p}) + e^{-p} \cdot \frac{p - (1 - e^{-p})}{e^{-p}} = p.
\]
Therefore $S := \sum_{i=1}^m X_i \sim \mathrm{Bin}(m,p)$, and $(S, N)$ is a coupling of $\mathrm{Bin}(m,p)$ and $\mathrm{Poi}(mp)$.

Under this coupling, $S \neq N$ can only happen if for some $i$ either (i) $N_i \geq 2$ (then $N$ counts at
least 2 from index $i$ but $S$ counts at most 1), or (ii) $N_i = 0$ but $X_i = 1$ (then $S$ counts 1 but $N$
counts 0). By the union bound,
\[
P(S \neq N) \leq m\bigl(P(N_1 \geq 2) + P(N_1 = 0, X_1 = 1)\bigr).
\]
Now $P(N_1 \geq 2) = 1 - e^{-p}(1 + p)$ and $P(N_1 = 0, X_1 = 1) = e^{-p} q = p - (1 - e^{-p})$ by definition of $q$.
Summing gives
\[
P(S \neq N) \leq mp(1 - e^{-p}).
\]
Apply Lemma~2.2 to conclude the first inequality in \eqref{eq:A1}. The second follows from $1 - e^{-p} \leq p$.
\end{proof}

\begin{lemma}[Poisson parameter perturbation]
Let $N \sim \mathrm{Poi}(\lambda)$ and $N' \sim \mathrm{Poi}(\lambda')$ with $\lambda, \lambda' \geq 0$. Then
\[
\mathrm{TV}(\mathrm{Poi}(\lambda), \mathrm{Poi}(\lambda')) \leq 1 - e^{-|\lambda - \lambda'|} \leq |\lambda - \lambda'|.
\]
\end{lemma}

\begin{proof}
Assume without loss of generality that $\lambda' \geq \lambda$. Let $M \sim \mathrm{Poi}(\lambda)$ and $R \sim \mathrm{Poi}(\lambda' - \lambda)$ be
independent, and set $M' := M + R$. Then $M' \sim \mathrm{Poi}(\lambda')$. Under this coupling, $M \neq M'$ iff $R \geq 1$,
hence
\[
P(M \neq M') = 1 - e^{-(\lambda' - \lambda)} = 1 - e^{-|\lambda - \lambda'|}.
\]
Apply Lemma~2.2 for the first inequality. The second uses $1 - e^{-x} \leq x$.
\end{proof}

\begin{lemma}[Multinomial $\to$ independent Poissons on rare categories]
Let $X \sim \mathrm{Mult}(m, p)$ on a finite alphabet $A$ with $p \in \Delta(A)$. Fix a subset $B \subseteq A$ and write $p_B := \sum_{b \in B} p_b$. Let $U = (U_b)_{b \in B}$ have independent coordinates $U_b \sim \mathrm{Poi}(m p_b)$. Then
\begin{equation}
\mathrm{TV}\bigl(\mathcal{L}((X_b)_{b \in B}), \mathcal{L}(U)\bigr) \leq m p_B (1 - e^{-p_B}). \label{eq:A3}
\end{equation}
\end{lemma}

\begin{proof}
Condition on the total number of draws that fall in $B$. Let $S := \sum_{b \in B} X_b$. Then $S \sim \mathrm{Bin}(m, p_B)$ and, given $S = s$, the vector $(X_b)_{b \in B}$ is multinomial $\mathrm{Mult}(s, p(\cdot|B))$ over $B$ with conditional probabilities $p_b/p_B$.

On the Poisson side, let $(U_b)_{b \in B}$ be independent with $U_b \sim \mathrm{Poi}(m p_b)$ and set $N := \sum_{b \in B} U_b$.
Then $N \sim \mathrm{Poi}(m p_B)$. Conditional on $N = s$, the vector $(U_b)_{b \in B}$ is also multinomial $\mathrm{Mult}(s, p(\cdot|B))$
(Poisson splitting).

Therefore, the only discrepancy between $(X_b)_{b \in B}$ and $(U_b)_{b \in B}$ comes from the discrepancy
between the totals $S$ and $N$; conditional on $S = N$ the conditional allocations match exactly.
Formally, couple $S$ and $N$ using Lemma~A.1 with $(m, p) = (m, p_B)$ so that $P(S \neq N) \leq m p_B(1 - e^{-p_B})$. Given $(S, N)$, sample the conditional multinomial allocation inside $B$ using the same auxiliary
randomness when $S = N$, and arbitrarily otherwise. This produces a coupling of $(X_b)_{b \in B}$ and
$(U_b)_{b \in B}$ whose mismatch probability is at most $P(S \neq N)$. Apply Lemma~2.2 to obtain \eqref{eq:A3}.
\end{proof}

\section{Privacy-curve convergence for the full hybrid experiment}
\label{app:hybrid-privacy}

All spaces in this appendix are finite or finite-dimensional Euclidean, hence standard Borel. In
particular, the regular conditional laws used below exist.

\begin{lemma}[Common independent factor does not affect privacy curves]
\label{lem:common-independent-factor}
Let $(P,Q)$ be a binary experiment and let $R$ be any probability measure independent of both $P$ and $Q$. Then for every $\varepsilon \ge 0$,
\[
\delta_{Q \otimes R \| P \otimes R}(\varepsilon) = \delta_{Q \| P}(\varepsilon).
\]
\end{lemma}

\begin{proof}
For any measurable $A\subseteq\mathcal X\times\mathcal Z$, write
\[
A_z:=\{x\in\mathcal X:(x,z)\in A\}.
\]
Then
\[
(Q\otimes R)(A)-e^\varepsilon(P\otimes R)(A)
=
\int\!\bigl(Q(A_z)-e^\varepsilon P(A_z)\bigr)\,R(dz)
\le \delta_{Q\|P}(\varepsilon)
\]
by the variational formula \eqref{eq:delta}. Taking the supremum over $A$ gives
\[
\delta_{Q\otimes R\|P\otimes R}(\varepsilon)\le \delta_{Q\|P}(\varepsilon).
\]
For the reverse inequality, fix any measurable $B\subseteq\mathcal X$ and take
$A=B\times\mathcal Z$. Then
\[
(Q\otimes R)(A)-e^\varepsilon(P\otimes R)(A)=Q(B)-e^\varepsilon P(B).
\]
Taking the supremum over $B$ yields
\[
\delta_{Q\otimes R\|P\otimes R}(\varepsilon)\ge \delta_{Q\|P}(\varepsilon).
\]
The two inequalities prove the claim.
\end{proof}

We keep the notation of Proposition~\ref{prop:hybrid-gaussian-compound-poisson}. In particular,
\[
G_n:=n^{-1/2}\Pi_G\widehat H_{n,k_n},
\qquad
Z_n:=\Pi_J\widehat H_{n,k_n},
\qquad
S_n=(G_n,Z_n),
\]
and under the neighboring alternative $T_{n,k_n+1}$ we use the same centering by $(n-k_n)\mu_0+k_n\mu_1$. We also write
\[
P_n^J:=\mathcal L(Z_n)\ \text{under }T_{n,k_n},
\qquad
Q_n^J:=\mathcal L(Z_n)\ \text{under }T_{n,k_n+1}.
\]
The next lemma resolves the interior-composition regime $\pi\in(0,1)$; the boundary cases are discussed separately below.

\begin{lemma}[Interior conditional smoothing for the dominant block]
\label{lem:interior-conditional-smoothing}
Assume the setting of Proposition~\ref{prop:hybrid-gaussian-compound-poisson} and, in addition, $\pi_n\to\pi\in(0,1)$. Define
\[
\theta_{b,n}:=\frac{W_b^{(n)}(y_{ba})}{W_b^{(n)}(y_{ba})+W_b^{(n)}(y_{bb})},
\qquad b\in\{0,1\}.
\]
For each $z\in M^\perp$ in the support of $Z_n$, let
\[
T_{b,n}(z):=N_{n,k_n}(D_b),
\qquad b\in\{0,1\},
\]
noting that $T_{b,n}(z)$ is $\sigma(Z_n)$-measurable because $\ker \Pi_J=M=\mathrm{span}\{g_0,g_1\}$ and adding an element of $M$ changes only the within-pair differences on $D_0,D_1$, not the pair totals. Let $R_{n,z}$ be the law of
\[
\Gamma_{n,z}:=\Psi_{n,z}(X_{0,n}^\circ,X_{1,n}^\circ),
\]
where $X_{0,n}^\circ\sim \mathrm{Bin}(T_{0,n}(z),\theta_{0,n})$ and $X_{1,n}^\circ\sim \mathrm{Bin}(T_{1,n}(z),\theta_{1,n})$ are independent and
\[
\Psi_{n,z}(x_0,x_1):=n^{-1/2}\Bigl((x_0-p_0T_{0,n}(z))g_0+(x_1-p_1T_{1,n}(z))g_1\Bigr).
\]
Then there exists a finite constant $C_{\mathrm{int}}=C_{\mathrm{int}}(p_0,p_1,\pi,\{\alpha_b(y)\})$ such that, for all sufficiently large $n$,
\begin{equation}
\int \mathrm{TV}\bigl(Q_{n,z}^{G\mid Z},R_{n,z}\bigr)\,Q_n^J(dz)
+\int \mathrm{TV}\bigl(P_{n,z}^{G\mid Z},R_{n,z}\bigr)\,P_n^J(dz)
\le \frac{C_{\mathrm{int}}}{\sqrt n},
\label{eq:B3-main}
\end{equation}
where $P_{n,z}^{G\mid Z}$ and $Q_{n,z}^{G\mid Z}$ are regular conditional laws of $G_n$ given $Z_n=z$ under $T_{n,k_n}$ and $T_{n,k_n+1}$, respectively.
\end{lemma}

\begin{proof}
Fix one of the two hypotheses and let $m_{0,n}^\star,m_{1,n}^\star$ denote the corresponding group sizes:
\[
(m_{0,n}^\star,m_{1,n}^\star)=
\begin{cases}
(m_{0,n},m_{1,n}), & \text{under }T_{n,k_n},\\
(m_{0,n}-1,m_{1,n}+1), & \text{under }T_{n,k_n+1}.
\end{cases}
\]
Let $L_{0,n}$ be the number of $0$-input users whose output falls in $\mathcal Y\setminus D_0$, and let $L_{1,n}$ be the number of $1$-input users whose output falls in $\mathcal Y\setminus D_1$. Let $A_n$ be the number of $1$-input users whose output falls in $D_0$, and let $B_n$ be the number of $0$-input users whose output falls in $D_1$. Then $A_n\le L_{1,n}$ and $B_n\le L_{0,n}$. Conditioning on the full rare configuration (equivalently: on $Z_n=z$ together with the latent assignment of rare outputs to the two input groups), the counts
\[
X_{0,n}:=N_{n,k_n}(y_{0a}),
\qquad
X_{1,n}:=N_{n,k_n}(y_{1a})
\]
are independent and satisfy
\begin{align*}
X_{0,n}
&\overset d= \mathrm{Bin}(m_{0,n}^\star-L_{0,n},\theta_{0,n})+\mathrm{Bin}(A_n,\vartheta_{10,n}),\\
X_{1,n}
&\overset d= \mathrm{Bin}(m_{1,n}^\star-L_{1,n},\theta_{1,n})+\mathrm{Bin}(B_n,\vartheta_{01,n}),
\end{align*}
where
\[
\vartheta_{10,n}:=\frac{W_1^{(n)}(y_{0a})}{W_1^{(n)}(y_{0a})+W_1^{(n)}(y_{0b})},
\qquad
\vartheta_{01,n}:=\frac{W_0^{(n)}(y_{1a})}{W_0^{(n)}(y_{1a})+W_0^{(n)}(y_{1b})},
\]
with an arbitrary value when the denominator is zero (then the corresponding count $A_n$ or $B_n$ is automatically zero). Moreover,
\[
T_{0,n}(z)=(m_{0,n}^\star-L_{0,n})+A_n,
\qquad
T_{1,n}(z)=(m_{1,n}^\star-L_{1,n})+B_n.
\]
Hence $R_{n,z}$ is exactly the law obtained by replacing the cross terms $\mathrm{Bin}(A_n,\vartheta_{10,n})$ and $\mathrm{Bin}(B_n,\vartheta_{01,n})$ by the native terms $\mathrm{Bin}(A_n,\theta_{0,n})$ and $\mathrm{Bin}(B_n,\theta_{1,n})$.

Write
\[
S_{0,n}\sim \mathrm{Bin}(m_{0,n}^\star-L_{0,n},\theta_{0,n}),
\qquad
V_{0,n}\sim \mathrm{Bin}(A_n,\vartheta_{10,n}),
\qquad
V_{0,n}^\circ\sim \mathrm{Bin}(A_n,\theta_{0,n}),
\]
independently, and analogously define $S_{1,n},V_{1,n},V_{1,n}^\circ$ on the $D_1$-pair. Since
\[
S_{0,n}+V_{0,n}^\circ\sim \mathrm{Bin}(T_{0,n}(z),\theta_{0,n}),
\qquad
S_{1,n}+V_{1,n}^\circ\sim \mathrm{Bin}(T_{1,n}(z),\theta_{1,n}),
\]
it suffices, by contraction under the measurable map $\Psi_{n,z}$ and by Lemma~\ref{lem:tv-tensorization}, to control the two one-dimensional distances
\[
\mathrm{TV}\bigl(\mathcal L(S_{0,n}+V_{0,n}),\mathcal L(S_{0,n}+V_{0,n}^\circ)\bigr),
\qquad
\mathrm{TV}\bigl(\mathcal L(S_{1,n}+V_{1,n}),\mathcal L(S_{1,n}+V_{1,n}^\circ)\bigr).
\]
For the first term, couple $V_{0,n}$ and $V_{0,n}^\circ$ as sums of $A_n$ coupled Bernoulli pairs, one pair for each cross message, so that
\[
E\bigl|V_{0,n}-V_{0,n}^\circ\mid Z_n,\text{rare configuration}\bigr|
\le A_n\,|\vartheta_{10,n}-\theta_{0,n}|
\le A_n.
\]
Conditional on the coupled values $(V_{0,n},V_{0,n}^\circ)=(v,w)$, repeated one-step shifts give
\[
\mathrm{TV}\bigl(\mathcal L(S_{0,n}+v),\mathcal L(S_{0,n}+w)\bigr)
\le |v-w|\,\sup_k P(S_{0,n}=k).
\]
Taking conditional expectation yields
\[
\mathrm{TV}\bigl(\mathcal L(S_{0,n}+V_{0,n}),\mathcal L(S_{0,n}+V_{0,n}^\circ)\bigr)
\le A_n\,\sup_k P(S_{0,n}=k).
\]
The same argument on $D_1$ gives
\[
\mathrm{TV}\bigl(\mathcal L(S_{1,n}+V_{1,n}),\mathcal L(S_{1,n}+V_{1,n}^\circ)\bigr)
\le B_n\,\sup_k P(S_{1,n}=k).
\]
Now $\theta_{b,n}\to p_b\in(0,1)$, so for all large $n$,
\[
\theta_{b,n}(1-\theta_{b,n})\ge \frac12 p_b(1-p_b),
\qquad b\in\{0,1\},
\]
and the unimodal binomial mass bound gives
\[
\sup_k P(S_{b,n}=k)
\le \frac{c_b}{\sqrt{m_{b,n}^\star-L_{b,n}+1}},
\qquad
c_b:=\sqrt{\frac{2}{p_b(1-p_b)}}.
\]
Therefore, conditional on the rare configuration,
\begin{equation}
\mathrm{TV}\bigl(\mathcal L(G_n\mid Z_n,\text{rare configuration}),R_{n,Z_n}\bigr)
\le c_0\frac{A_n}{\sqrt{m_{0,n}^\star-L_{0,n}+1}}+c_1\frac{B_n}{\sqrt{m_{1,n}^\star-L_{1,n}+1}}.
\label{eq:B3-conditional-bound}
\end{equation}
Averaging over the latent rare configuration and then over $Z_n$ (tower property) gives the same bound for the conditional laws $P_{n,z}^{G\mid Z}$ or $Q_{n,z}^{G\mid Z}$.

It remains to integrate the right-hand side. Set
\[
\kappa:=\frac14\min\{\pi,1-\pi\}>0.
\]
Since $\pi_n\to\pi$, for all sufficiently large $n$ we have $m_{0,n}^\star\ge 3\kappa n$ and $m_{1,n}^\star\ge 3\kappa n$ under both hypotheses. Let $R_n:=L_{0,n}+L_{1,n}$. On the event $\{R_n\le \kappa n\}$,
\[
\frac{A_n}{\sqrt{m_{0,n}^\star-L_{0,n}+1}}+\frac{B_n}{\sqrt{m_{1,n}^\star-L_{1,n}+1}}
\le \frac{A_n+B_n}{\sqrt{2\kappa n}}
\le \frac{R_n}{\sqrt{2\kappa n}}.
\]
On the complementary event, the same quantity is bounded by $R_n$. Hence
\[
E\Bigl[\frac{A_n}{\sqrt{m_{0,n}^\star-L_{0,n}+1}}+\frac{B_n}{\sqrt{m_{1,n}^\star-L_{1,n}+1}}\Bigr]
\le \frac{E[R_n]}{\sqrt{2\kappa n}}+E\bigl[R_n\mathbf 1_{\{R_n>\kappa n\}}\bigr].
\]
Since $R_n$ is a Poisson-binomial sum of rare-output indicators,
\[
E[R_n]\le 2(\Lambda_0+\Lambda_1),
\qquad
E[R_n^2]\le 2(\Lambda_0+\Lambda_1)+4(\Lambda_0+\Lambda_1)^2
\]
for all large $n$, where
\[
\Lambda_b:=\sum_{y\notin D_b}\alpha_b(y), \qquad b\in\{0,1\}.
\]
By Cauchy--Schwarz and Markov,
\[
E\bigl[R_n\mathbf 1_{\{R_n>\kappa n\}}\bigr]
\le \frac{E[R_n^2]}{\kappa n}
=O(n^{-1}).
\]
Combining this with \eqref{eq:B3-conditional-bound} gives \eqref{eq:B3-main}. One may take, for example,
\[
C_{\mathrm{int}}
:=\Bigl(\sqrt{\frac{2}{p_0(1-p_0)}}+\sqrt{\frac{2}{p_1(1-p_1)}}\Bigr)
\Bigl(\frac{2(\Lambda_0+\Lambda_1)}{\sqrt{2\kappa}}+\frac{2(\Lambda_0+\Lambda_1)+4(\Lambda_0+\Lambda_1)^2}{\kappa}\Bigr).
\]
Since $\mathcal Y$ is finite, this depends only on $p_0,p_1,\pi$ and finitely many $\alpha_b(y)$; in particular it is controlled by $p_0,p_1,\pi$ and $\sup_{b,y}\alpha_b(y)$ once $\mathcal Y$ is fixed.
\end{proof}


\begin{remark}[Boundary compositions $\pi\in\{0,1\}$]
\label{rem:B3-boundary}
Lemma~\ref{lem:interior-conditional-smoothing} is an interior statement and is false at the boundary. Indeed, at $\pi=0$ one can choose rare transitions $0\to D_1$ with two positive intensities $\alpha_0(y_{1a}),\alpha_0(y_{1b})$ and choose the native split parameter $p_1$ so that
\[
p_1\neq \frac{\alpha_0(y_{1a})}{\alpha_0(y_{1a})+\alpha_0(y_{1b})}.
\]
On the projected jump event $Z_n=\Delta$, the unique $D_1$-message is then generated under $Q_n$ by the distinguished $1$-user and hence splits asymptotically with parameter $p_1$, whereas under $P_n$ it is generated by a rare $0\to D_1$ error and hence splits with the rare-error parameter above. Thus the conditional total variation stays bounded away from $0$. This causes no problem for privacy curves, however, because the boundary compositions are already covered by Sections~3--4: part~(b) of Corollary~\ref{cor:multi-curve-series}, together with Remark~\ref{rem:boundary-scalar-privacy}, identifies the boundary limit experiment as a scalar Poisson-shift on the switched coordinate plus common independent noise on the remaining coordinates. Hence the boundary privacy curve is exactly the one treated in Theorem~\ref{thm:poisson-shift}; in the binary case this is also the boundary specialization of Theorem~\ref{thm:skellam-shift}.
\end{remark}

\begin{corollary}[Privacy-curve convergence for the full hybrid experiment]
\label{cor:full-hybrid-privacy-convergence}
Assume the setting of Proposition~\ref{prop:hybrid-gaussian-compound-poisson}. Then for every fixed $\varepsilon\ge 0$,
\[
\delta_{Q_n\|P_n}(\varepsilon)\longrightarrow \delta_{Q_\infty\|P_\infty}(\varepsilon),
\qquad
P_\infty=\mathcal L(G,J),\quad Q_\infty=\mathcal L(G,J+\Delta).
\]
Equivalently,
\[
\delta_{Q_n\|P_n}(\varepsilon)\longrightarrow \delta_{\mathcal L(J+\Delta)\|\mathcal L(J)}(\varepsilon).
\]
More precisely:
\begin{enumerate}
\item[(i)] if $\pi\in(0,1)$, then for all sufficiently large $n$,
\[
\bigl|\delta_{Q_n\|P_n}(\varepsilon)-\delta_{Q_n^J\|P_n^J}(\varepsilon)\bigr|
\le \frac{C_{\mathrm{int}}(1+e^\varepsilon)}{\sqrt n},
\]
with $C_{\mathrm{int}}$ from Lemma~\ref{lem:interior-conditional-smoothing};
\item[(ii)] if $\pi\in\{0,1\}$, then the full privacy curve reduces to the boundary Poisson-shift case already covered in Sections~3--4.
\end{enumerate}
\end{corollary}

\begin{proof}
For $\pi\in(0,1)$, define binary experiments with the common conditional factor $R_{n,z}$ by
\[
\widetilde P_n(dg,dz):=P_n^J(dz)R_{n,z}(dg),
\qquad
\widetilde Q_n(dg,dz):=Q_n^J(dz)R_{n,z}(dg).
\]
Let $K((g,z),B):=\mathbf 1_B(z)$ be the projection kernel from $M\times M^\perp$ onto $M^\perp$, and let
\[
L(z,A):=\int \mathbf 1_A(g,z)\,R_{n,z}(dg)
\]
be the enrichment kernel from $M^\perp$ back to $M\times M^\perp$. By construction,
\[
\widetilde P_n K=P_n^J,\qquad \widetilde Q_n K=Q_n^J,\qquad P_n^J L=\widetilde P_n,\qquad Q_n^J L=\widetilde Q_n.
\]
Hence data processing in the two directions gives
\[
\delta_{\widetilde Q_n\|\widetilde P_n}(\varepsilon)=\delta_{Q_n^J\|P_n^J}(\varepsilon).
\]
By Lemma~\ref{lem:privacy-curve-stability},
\[
\bigl|\delta_{Q_n\|P_n}(\varepsilon)-\delta_{\widetilde Q_n\|\widetilde P_n}(\varepsilon)\bigr|
\le \mathrm{TV}(Q_n,\widetilde Q_n)+e^\varepsilon\mathrm{TV}(P_n,\widetilde P_n).
\]
The two TV terms are exactly the two integrals in \eqref{eq:B3-main}, so Lemma~\ref{lem:interior-conditional-smoothing} yields
\[
\bigl|\delta_{Q_n\|P_n}(\varepsilon)-\delta_{Q_n^J\|P_n^J}(\varepsilon)\bigr|
\le \frac{C_{\mathrm{int}}(1+e^\varepsilon)}{\sqrt n}.
\]
Now Corollary~\ref{cor:projected-jump-lecam} gives $P_n^J\to P_\infty^J$ and $Q_n^J\to Q_\infty^J$ in total variation, where
\[
P_\infty^J=\mathcal L(J),
\qquad
Q_\infty^J=\mathcal L(J+\Delta).
\]
Applying Lemma~\ref{lem:privacy-curve-stability} to the projected experiment shows
\[
\delta_{Q_n^J\|P_n^J}(\varepsilon)\to \delta_{Q_\infty^J\|P_\infty^J}(\varepsilon).
\]
Finally, Lemma~\ref{lem:common-independent-factor} removes the common Gaussian factor in the limit:
\[
\delta_{Q_\infty\|P_\infty}(\varepsilon)=\delta_{\mathcal L(G,J+\Delta)\|\mathcal L(G,J)}(\varepsilon)
=\delta_{\mathcal L(J+\Delta)\|\mathcal L(J)}(\varepsilon).
\]
This proves the interior case.

If $\pi\in\{0,1\}$, Remark~\ref{rem:B3-boundary} identifies the boundary privacy problem with the scalar Poisson-shift setting already treated in Theorem~\ref{thm:poisson-shift}; equivalently, in the binary case it is the boundary specialization of Theorem~\ref{thm:skellam-shift}. Hence the same convergence conclusion holds at the boundary as well.
\end{proof}

\begin{thebibliography}{99}

\bibitem{Shv26} A.~Shvets. Universal Shuffle Asymptotics, Part I: Sharp Privacy Analysis in the Gaussian Regime. \textit{arXiv:2602.09029}, 2026.

\bibitem{CGS11} L.~H.~Y.~Chen, L.~Goldstein, and Q.-M.~Shao. \textit{Normal Approximation by Stein's Method}. Springer, 2011.

\bibitem{LeC86} L.~Le~Cam. \textit{Asymptotic Methods in Statistical Decision Theory}. Springer, 1986.

\bibitem{BHJ92} A.~D.~Barbour, L.~Holst, and S.~Janson. \textit{Poisson Approximation}. Oxford University Press, 1992.

\bibitem{DRS22} J.~Dong, A.~Roth, and W.~J.~Su. Gaussian differential privacy. \textit{Journal of the Royal Statistical Society: Series B}, 84(1):3--37, 2022.

\bibitem{FMT21} V.~Feldman, A.~McMillan, and K.~Talwar. Hiding among the clones: A simple and nearly optimal analysis of privacy amplification by shuffling. In \textit{FOCS}, 2021.

\bibitem{EFM19} U.~Erlingsson, V.~Feldman, I.~Mironov, A.~Raghunathan, K.~Talwar, and A.~Thakurta. Amplification by shuffling: From local to central differential privacy via anonymity. In \textit{SODA}, 2019.

\bibitem{CSU19} A.~Cheu, A.~Smith, J.~Ullman, D.~Zeber, and M.~Zhilyaev. Distributed differential privacy via shuffling. In \textit{EUROCRYPT}, 2019.

\bibitem{BBG19} B.~Balle, J.~Bell, A.~Gascon, and K.~Nissim. The privacy blanket of the shuffle model. In \textit{CRYPTO}, 2019.

\bibitem{GKL21} A.~Girgis, N.~Kairouz, Z.~Liu, T.~Steinke, and K.~Talwar. On the R\'{e}nyi differential privacy of the shuffle model. In \textit{ACM CCS}, 2021.

\bibitem{Mir17} I.~Mironov. R\'{e}nyi differential privacy. In \textit{CSF}, 2017.

\bibitem{Pet75} V.~V.~Petrov. \textit{Sums of Independent Random Variables}. Springer, 1975.

\bibitem{vdV98} A.~W.~van der~Vaart. \textit{Asymptotic Statistics}. Cambridge University Press, 1998.

\bibitem{Rol07} A.~R\"{o}llin. Translated Poisson approximation using exchangeable pair couplings. \textit{Annals of Applied Probability}, 17(5/6):1596--1614, 2007.

\bibitem{TL26} S.~Takagi and S.~P.~Liew. Analysis of shuffling beyond pure local differential privacy. \textit{arXiv:2601.19154}, 27~January~2026.

\end{thebibliography}

\end{document}
